\documentclass{article}
% === Overleaf에서 컴파일 엔진을 XeLaTeX으로 설정 필요 ===

% --- 문서 크기 설정 ---
\usepackage[a4paper,left=1.5cm,right=1.5cm,top=2cm,bottom=2cm]{geometry}

% --- 한글 / 영문 폰트 (kotex 사용 - 한글 처리 최적화) ---
\usepackage{kotex}
\setmainfont{Noto Serif}
\setmainhangulfont{Noto Serif CJK KR}
\setsansfont{Noto Sans}
\setsanshangulfont{Noto Sans CJK KR}
\setmonofont{Noto Sans Mono}
\setmonohangulfont{Noto Sans Mono CJK KR}

% --- 한글 줄바꿈 및 띄어쓰기 개선 ---
\XeTeXlinebreaklocale "ko_KR"
\XeTeXlinebreakskip = 0pt plus 1pt minus 0.1pt
\XeTeXlinebreakpenalty = 0

% --- 미세 타이포그래피 조정 ---
\usepackage{microtype}

% --- 수학 및 양자 관련 ---
\usepackage{amsmath, amssymb, amsfonts, amsthm, bm, nicefrac}
\usepackage{braket}

% --- 표/정렬/목차 등 ---
\usepackage{multirow}
\usepackage{tabularx}
\usepackage{array}
\usepackage{booktabs}
\usepackage{enumitem}
\setlist[itemize,2]{label=\(\circ\)} % 작은 빈 원


% --- 그래픽 / 색 / TikZ ---
\usepackage{graphicx}
\usepackage{xcolor}
\usepackage{soul}
\usepackage[export]{adjustbox}
\usepackage{tikz}
\usetikzlibrary{
    positioning,
    arrows.meta,
    fit,
    backgrounds,
    calc,
    automata,
    shapes.geometric,
    decorations.pathmorphing,
    decorations.pathreplacing
}
\pgfdeclarelayer{background}
\pgfsetlayers{background,main}
\usepackage{float}  % figure 위치 제어 [H]
\usepackage{caption}
\usetikzlibrary{positioning}

% --- 인용 (hyperref보다 먼저 로드) --- 
\usepackage[backend=biber, style=numeric, sorting=nyt]{biblatex}
\addbibresource{Setting/citation.bib}

% --- 하이퍼링크 (보통 맨 마지막에 로드) ---
\usepackage{hyperref}
\hypersetup{
    colorlinks=true,
    linkcolor=blue,
    filecolor=magenta,
    urlcolor=blue,
    pdftitle={Military Project},
    pdfpagemode=FullScreen,
}

% --- 색 정의 ---
\definecolor{SkyBlue}{rgb}{0.5, 0.8, 1}
\definecolor{MyPink}{rgb}{1, 0.682, 0.843}
\definecolor{customblue}{HTML}{1f77b4}
\definecolor{customorange}{HTML}{ff7f0e}
\definecolor{customgreen}{HTML}{2ca02c}
\definecolor{custompurple}{HTML}{C4A3F0}

% --- 섹션 번호 깊이 ---
\setcounter{secnumdepth}{2}

% --- 정리된 theorem 등 정의 ---
\newtheorem{theorem}{Theorem}[section]
\newtheorem{lemma}[theorem]{Lemma}
\newtheorem{definition}{Definition}[section]
\newtheorem{prop}{Proposition}
\newtheorem{axiom}{Axiom}
\newtheorem{claim}{Claim}
\newtheorem{corollary}{Corollary}[theorem]

% --- 커맨드 ---
\newcommand{\Theme}[1]{{\vspace{0.5cm}\noindent\textbf{#1}}}


\title{특허 명세서: 골전도 음향 장치 및 딥러닝을 활용한 극단적 소음 제거 방법 및 이를 이용한 음성 통신 장치}
\author{일병 박성민 \\ 육군본부 직할 미래혁신연구센터 군사과학기술연구대}
\date{January 19, 2026 -- \today}


\begin{document}
\maketitle
\tableofcontents





\newpage
\part{서지사항}
\section{발명의 명칭}
\begin{itemize}
    \item \textbf{국문}: 골전도 음향 장치 및 딥러닝을 활용한 극단적 소음 제거 방법 및 이를 이용한 음성 통신 장치
    \item \textbf{영문}: Denoising Method for Extreme-Noise Environments Using a Bone-Conduction Acoustic Device and Deep Learning Model, and Voice Communication Device Thereof
\end{itemize}






\part{발명의 상세한 설명}


\section{기술분야}

본 발명(골전도 음향 장치 및 딥러닝을 활용한 극단적 소음 제거 방법 및 이를 이용한 음성 통신 장치)은 극한 소음 환경에서의 음성 신호 처리 기술 및 음성 통신 장치에 관한 것이다.
구체적으로, 비선형 왜곡 및 마이크·스피커의 포화가 발생하는 극한 소음 환경에서 사용자의 청력을 보호하면서, 소음 내성 음향 장치부(골전도 장치, 지향성 마이크로폰 및 마이크 어레이 등)와 딥러닝 기반 음성 향상 모델을 이용한 소음 제거 방법 및 이를 적용한 음성 통신 장치에 관한 것이다.
% [수정] 지적 #1 대응: "생성적 음성 복원 모델" → "음성 향상 모델" 용어 통일

\section{발명의 배경이 되는 기술}

\subsection{발명의 배경} 
극단적 소음 환경은 통상적인 생활 소음 환경에서 나타나지 않는 고유한 물리적·기술적 문제를 수반한다.
본 발명은 이러한 문제를 해결하고자 하는 배경에서 고안되었다.



\paragraph{1. 사용자의 청력 보호와 의사소통 단절 딜레마}

현대전의 전장, 항공 우주 발사체 주변, 중공업 현장은 120 dB 이상의 극심한 소음이 상시 발생하는 환경이다. 
이러한 고강도 소음은 인체에 즉각적이고 누적적인 청력 손상을 유발할 위험이 크므로, 작업자나 전투원은 귀마개 또는 귀덮개와 같은 청력 보호 장비를 필수로 착용해야 한다. 
이는 장기적으로 청력 보존뿐만 아니라, 위험 요소가 공존하는 극한 환경에서의 생존성 확보 및 업무(전투) 수행 능력을 유지하기 위한 필수적인 요건이다.

그러나, 대부분의 종래 청력 보호 장비는 귀를 물리적으로 폐쇄하는 수동적 차음 방식(passive noise isolation)에 주로 의존하고 있어, 소음뿐만 아니라 음성을 포함한 모든 대역의 소리를 감쇠시킨다는 근본적인 한계가 있다. 이로 인해 사용자의 청력은 보호될 수 있으나, 동료 간의 명확한 의사소통이나 주변의 경고음 인지가 어려워지는 문제가 발생한다. 이는 결과적으로 업무 효율을 저하시킬 뿐만 아니라, 긴박한 상황에서 지휘 통제의 지연이나 상황 판단의 오류를 초래하여 인명 사고나 작전 실패로 이어질 위험성을 증대시킨다.

또한, 종래의 음성 통신 기술은 주로 저소음 또는 통제된 소음 환경에 최적화되어 있어, 상기한 극한 소음 환경에서 청력 보호와 고품질의 음성 전달을 동시에 달성하는 데 어려움이 있다. 
일부 장비가 음성 증폭이나 필터링 기능을 제공하고 있으나, 물리적 보호 장비 착용에 따른 음성 수용력 저하 문제를 근본적으로 해결하지 못하며, 고소음 유입 시 음성 명료도(speech intelligibility)가 현저히 떨어지는 실정이다. 아울러, 귀에 직접 착용하는 방식은 헬멧 등 타 보호 장비와의 호환성 문제를 야기하거나 장시간 착용 시 사용자에게 불편함을 주는 문제점이 있다.



\paragraph{2. 극한 소음의 비선형·포화 특성에 의한 기존 소음 제거 기술 한계}

\begin{table}[b]
    \centering
    \renewcommand{\arraystretch}{1.3} % 행 간격 약간 넓힘
    \begin{tabular}{c c c l} % 세로선 제거, 깔끔한 스타일
        \hline
        \textbf{환경 구분} & \textbf{주요 소음원 (예시)} & \textbf{대표 음압 범위 (dB)}& \textbf{주요 물리적/신호적 특성} \\
        \hline
        \hline
        \textbf{전장} & 총성, 포격, 폭발음 & $140 \sim 170^+$ & 충격음의 비선형성, 임펄스 특성  \\
        \hline
        \textbf{우주산업} & 발사체, 제트, 소닉붐 & $130 \sim 160$ & 광대역 난류 소음, 비선형 전파 \\
        \hline
        \textbf{중공업} & 해머, 프레스 가공, 발파 & $110 \sim 140$ & 충격성 잡음, 센서 동적범위 초과 \\
        \hline
    \end{tabular}
    \caption{극한 소음 환경의 예시 및 특성. (대표 음압 범위는 NIOSH, OSHA 등의 산업안전 문헌 및 군사 규격에서 보고된 대표적 수치를 참고하였다.)}
    \label{tab:extreme_noise_examples}
\end{table}

상기 표 [Tab.~\ref{tab:extreme_noise_examples}]에 예시된 극한 소음 환경은 통상적인 소음 환경과 물리적 특성이 확연히 구분되므로, 이에 대응하기 위한 새로운 기술적 접근이 요구된다. 
그러나 종래의 음향학 및 이에 기반한 통상적인 소음 제거 기술들은 음압 섭동 $p$가 충분히 작아 매질이 선형 탄성체로 거동한다는 가정에 기초한다. 
이는 아래의 식~\eqref{eq: linear wave eq.}과 같이 달랑베르 연산자(d'Alembertian, $\square$)를 이용한 선형 파동 방정식으로 기술된다.
(여기서 $c_{0}$는 평형상태의 음속, $\nabla^2$는 라플라시안 연산자를 의미한다.)
\begin{equation}\label{eq: linear wave eq.}
    \square p (\mathbf{x}, t) = 0, \qquad 
    \text{ where } \;
    \square \equiv \frac{1}{c_0^2}\frac{\partial^2}{\partial t^2} - \nabla^{2}
\end{equation}
위 식~\eqref{eq: linear wave eq.}이 성립하는 선형 체계 하에서는 중첩의 원리(Principle of Superposition)가 적용되므로, 특정 수신점에서의 관측 음압 $p$를 목적 음성 신호($s$)와 배경 잡음($n$)의 선형 결합으로 모델링할 수 있다.
\begin{equation}\label{eq: additive model}
    p (\mathbf{x}, t) = s(\mathbf{x}, t) + n (\mathbf{x}, t)
\end{equation}
위너 필터(Wiener filter), 스펙트럼 차감법(spectral subtraction) 등 고전적 최적 필터링 기법은 위 식~\eqref{eq: additive model}에 기초하여, 전체 신호에서 잡음 성분을 추정·제거하는 방식에 의존해 왔다. 이러한 기법들은 다음의 가정들을 전제로 한다.
\begin{enumerate}
    \item \textbf{선형성(Linearity)}: 음압 섭동 $p$가 대기압 $p_0$에 비해 충분히 작은 \textbf{미소 진폭} 조건을 만족하며, 매질이 선형 탄성체로 거동하여 식~\eqref{eq: linear wave eq.}에 따라 기술된다.
    한편, 상기 미소 진폭 조건이 성립하지 않는 강한 음압 수준에서의 음향 전파는 미소 섭동 이론을 벗어나 비선형 효과를 포함한 Westervelt 방정식으로 기술된다.
    \begin{equation}\label{eq: westervelt}
        \square p = -\frac{\delta}{c_0^4}\frac{\partial^3 p}{\partial t^3} - \frac{\beta}{\rho_0 c_0^4}\frac{\partial^2 p^2}{\partial t^2}
    \end{equation}
    여기서 $\delta$는 음향 확산계수, $\beta$는 매질의 비선형 계수, $\rho_{0}$는 평형상태의 밀도이다. 
    식~\eqref{eq: westervelt}의 우변에 나타나는 열점성 소산항(제1항)과 비선형 항(제2항)은 전파 과정에서 에너지 손실, 고조파(harmonics) 생성 및 파형의 비가역적 왜곡을 유발하며, 상기 선형성 가정을 근본적으로 위배한다.
    
    \item \textbf{가우시안 정상성(Gaussian Stationarity)}: 배경 잡음 $n(t)$가 정상(stationary) 가우시안 확률 과정을 따르며, 그 통계적 특성(평균, 분산, 스펙트럼)이 시간에 따라 일정하다.
    
    \item \textbf{협대역 잡음(Narrowband Noise)}: 잡음의 스펙트럼 에너지가 제한된 주파수 대역에 집중되어, 주파수 선택적 필터링이 유효하다.
    
    \item \textbf{센서 선형 응답(Sensor Linearity)}: 음향 센서가 
    동작 범위 내에서 입력 음압에 비례하는 출력을 생성하며, 
    포화(saturation) 또는 클리핑(clipping)이 발생하지 않는다.
    \item \textbf{시불변 채널(Time-Invariant Channel)}: 음향 전달 경로의 
    임펄스 응답이 시간에 따라 일정하여, 잔향 등 채널 효과의 추정 및 
    역산(deconvolution)이 가능하다.
\end{enumerate}
그러나, 앞선 표~\ref{tab:extreme_noise_examples}에서 제시한 바와 같이 음압 레벨(SPL)이 130dB를 초과하는 극한 소음 환경에서는 소음환경의 비선형성, 비정상성 등으로 인해 앞서 언급한 다섯 가지 가정들이 복합적으로 위배 된다.
이 경우, 해당 가정들의 성립을 전제로 설계된 고전적 선형 필터링 기법은 물리적 타당성과 이론적 최적성을 상실하게 되며, 이는 실질적인 음성 복원 성능의 현저한 저하로 이어진다.






\paragraph{3. 보호 장구 일체형 음향장비의 하드웨어적 한계}

청력 보호를 위해 이도를 폐쇄하거나 귀 전체를 덮는 보호 장구와, 통신을 위한 음향 입출력 장치(마이크로폰, 스피커)가 결합된 시스템은 120\,dB 이상의 극단적 소음 환경에서 다음과 같은 물리적·구조적 한계에 직면한다.

첫째, 입력부(마이크로폰, 이하 마이크)의 경우, 센서의 배치 위치 및 수음 방식에 따라 상이한 기술적 난제가 존재한다. 
우선, 외부 노출형 마이크는 고강도 음향 에너지에 진동판이 직접적으로 노출되므로, 센서의 허용 입력을 초과하는 \textbf{센서 포화(saturation)} 및 \textbf{클리핑(clipping) 현상}이 불가피하다. 
이러한 클리핑은 단순한 파형의 왜곡을 넘어, 파형의 피크 정보가 절단되는 비가역적 정보 손실을 초래한다. 
이는 후처리 필터링 알고리즘으로도 원본 신호의 복원이 불가능하므로, 긴박한 작전 상황에서 명령 전달 실패 등 치명적인 결과를 야기할 수 있다. 
또한, 총성이나 폭발음과 같은 고에너지 충격형 소음 유입 시, 센서 진동판이 과도 응답(transient response) 후 평형 상태로 복귀하기까지의 회복 시간 동안 음성 수집이 불가능한 수음 공백 구간이 발생할 수 있다.

한편, 폐쇄형 또는 접촉식 마이크는 외부 소음의 물리적 전달이 차단되어 센서 포화 문제는 완화되나, 신호 대역폭과 품질 측면에서 한계가 있다. 
특히, 사용자의 호흡, 심박, 저작(chewing) 등에 기인한 생체 잡음과, 장비와 신체 간의 접촉에 의한 마찰음이 음성 신호에 혼입되어 음성 명료도를 현저히 저하시키는 요인이 된다.

둘째, 출력부(스피커)의 경우, 통상의 공기 전도 방식은 청력 보호 장치에 의해 음향 전달 경로가 물리적으로 차단된다는 근본적 모순이 존재한다. 
이를 극복하기 위해 스피커 출력을 과도하게 높일 경우, 밀폐된 공간 내의 음압 상승으로 인한 청력 손상 위험이 발생하여 '\textbf{청력 보호}'라는 본연의 목적과 '\textbf{명료한 통신}'이라는 기능적 요구가 상충하는 한계가 있다.




\newpage
\section{선행 기술 문헌}

다음은 본 발명과 관련된 주요 선행기술 문헌의 목록이다.

\begin{table}[!h]
    \centering
    \renewcommand{\arraystretch}{1.2}
    \begin{tabular}{|c|p{6.5cm}|p{3.5cm}|p{3cm}|c|}
    \hline
    번호 & 발명의 명칭 (영문) & 출원인 & 공보번호 & 연도 \\
    \hline
    
    \multirow{3}{*}{1} 
    & \multirow{3}{=}{딥러닝 알고리즘을 이용한 소음 제거 방법 및 장치} 
    & \multirow{3}{=}{Mobilint Inc.} 
    & \href{https://patents.google.com/patent/KR102263135B1/en?oq=KR102263135B1}{KR102263135B1} & 2020 \\ \cline{4-5}
    &  &  & \href{https://patents.google.com/patent/WO2022124452A1/en?oq=KR102263135B1}{WO2022124452A1} & 2020 \\ \cline{4-5}
    &  &  & \href{https://patents.google.com/patent/US12412556B2/en?oq=KR102263135B1}{US12412556B2} & 2023 \\
    \hline
    
    2 
    & 골전도 진동자를 이용한 조종사용 헬멧 
    & 이준환 
    & \href{https://patents.google.com/patent/KR20110004763U/en?oq=KR20110004763U}{KR20110004763U} 
    & 2011 \\
    \hline
    
    \multirow{5}{*}{3} 
    & \multirow{5}{=}{골전도 센서 및 마이크로폰 신호를 융합한 딥러닝 음성 추출 및 노이즈 저감 방법} 
    & \multirow{5}{=}{Elevoc Technology Co., Ltd.} 
    & \href{https://patents.google.com/patent/EP4044181A4/en?oq=US20220392475A1}{EP4044181A1} & 2019 \\ \cline{4-5}
    &  &  & \href{https://patents.google.com/patent/WO2021068120A1/en?oq=US20220392475A1}{WO2021068120A1} & 2021 \\ \cline{4-5}
    &  &  & \href{https://patents.google.com/patent/US20220392475A1/en?oq=US20220392475A1}{US20220392475A1} & 2022 \\ \cline{4-5}
    &  &  & \href{https://patents.google.com/patent/KR102429152B1/en?oq=US20220392475A1}{KR102429152B1} & 2022 \\ \cline{4-5}
    &  &  & \href{https://patents.google.com/patent/JP2022505997A/en?oq=US20220392475A1}{JP2022505997A} & 2022 \\
    \hline
    
    \multirow{9}{*}{4} 
    & \multirow{9}{=}{헬멧을 위한 능동 소음 제거 시스템} 
    & \multirow{9}{=}{Daal Noise Control Systems AS} 
    & \href{https://patents.google.com/patent/CA3118977A1/en?oq=US12462783B2}{CA3118977A1} & 2020 \\ \cline{4-5}
    &  &  & \href{https://patents.google.com/patent/WO2020094733A1/en?oq=US12462783B2}{WO2020094733A1} & 2020 \\ \cline{4-5}
    &  &  & \href{https://patents.google.com/patent/GB2578744B/en?oq=US12462783B2}{GB2578744B} & 2022 \\ \cline{4-5}
    &  &  & \href{https://patents.google.com/patent/DK3876777T3/en?oq=US12462783B2}{DK3876777T3} & 2023 \\ \cline{4-5}
    &  &  & \href{https://patents.google.com/patent/EP3876777B8/en?oq=US12462783B2}{EP3876777B8} & 2023 \\ \cline{4-5}
    &  &  & \href{https://patents.google.com/patent/JP7422760B2/en?oq=US12462783B2}{JP7422760B2} & 2023 \\ \cline{4-5}
    &  &  & \href{https://patents.google.com/patent/KR102706589B1/en?oq=US12462783B2}{KR102706589B1} & 2024 \\ \cline{4-5}
    &  &  & \href{https://patents.google.com/patent/CN113260266B/en?oq=US12462783B2}{CN113260266B} & 2024 \\ \cline{4-5}
    &  &  & \href{https://patents.google.com/patent/US12462783B2/en?oq=US12462783B2}{US12462783B2} & 2025 \\ 
    \hline
    5 
    & 골전도 통신이 가능한 안전모 
    & 주식회사 오토스광학 
    & \href{https://patents.google.com/patent/KR200317070Y1/en?oq=KR200317070Y1}{KR200317070Y1} 
    & 2003 \\
    \hline
    
    % \multirow{3}{*}{6} 
    % & \multirow{3}{=}{골전도 마이크로폰이 내장된 헤드셋} 
    % & \multirow{3}{=}{Audio-Technica KK} 
    % & \href{https://patents.google.com/patent/US8325963B2/en?oq=US8325963B2}{US8325963B2} & 2012 \\ \cline{4-5}
    % &  &  & \href{https://patents.google.com/patent/US8325963B2/en?oq=JP5269618B2}{JP5269618B2} & 2013 \\ \cline{4-5}
    % &  &  & \href{https://patents.google.com/patent/CN101795143B/en?oq=JP5269618B2}{CN101795143B} & 2015 \\
    % \hline
    
    \multirow{5}{*}{6} 
    & \multirow{5}{=}{골전도 센서를 구비한 청각 기기} 
    & \multirow{5}{=}{GN Hearing A/S} 
    & \href{https://patents.google.com/patent/US20250210031A1/en?oq=EP3737115A1}{US20250210031A1} & 2025 \\ \cline{4-5}
    &  &  & \href{https://patents.google.com/patent/EP3967060A1/en?oq=EP3737115A1}{EP3737115A1} & 2020 \\ \cline{4-5}
    &  &  & \href{https://patents.google.com/patent/WO2020225294A1/en?oq=EP3737115A1}{WO2020225294A1} & 2020 \\ \cline{4-5}
    &  &  & \href{https://patents.google.com/patent/JP2022531363A/en?oq=EP3737115A1}{JP2022531363A} & 2022 \\ \cline{4-5}
    &  &  & \href{https://patents.google.com/patent/CN114009063A/en?oq=EP3737115A1}{CN114009063A} & 2022 \\
    \hline
    
    \end{tabular}
\caption{선행기술 특허 서지정보 요약}
\end{table}



\subsection{특허 문헌}

이하에서 분석하는 선행기술 문헌들은, 각각의 기술적 구성에서 차이가 있으나, 
본 발명이 해결하고자 하는 과제의 관점에서 다음과 같은 공통적 한계를 갖는다:
\begin{enumerate}[label = (\alph*)]
    \item 음압 레벨 130\,dB 이상의 극한 소음 환경에서의 동작을 구체적으로 고려하지 않으며,
    \item 센서 포화(saturation) 및 클리핑(clipping)에 의한 비가역적 파형 손실 상황에 대한 대응 수단을 개시하지 않고,
    
    \item 극한 소음 환경에서의 청력 보호, 골전도 양방향 음성 입출력, 센서 포화 대응, 및 딥러닝 기반 음성 향상을 하나의 통합된 체계로 결합하는 구성을 개시하지 않는다.
\end{enumerate}
이하에서는 각 문헌의 구체적인 기술 구성과 개별적 한계를 기술한다.

\subsubsection*{1. 딥러닝 알고리즘을 이용한 소음 제거 방법 및 장치 (KR 10-2263135 B1)}

대한민국 등록특허 제10-2263135호(이하 '135 특허)는 딥러닝 알고리즘을 이용하여 수집된 소음 신호로부터 음성 신호만을 추출한 제1 소리 신호와, 해당 신호 내 사람의 음성 신호가 포함되어 있을 확률 값 $P$를 산출하고, $P$ 값에 따라 제1 소리 신호 또는 로그 변환($y = \log(x + 1)$)된 제2 소리 신호를 선택적으로 출력하는 소음 제거 방법 및 장치를 개시한다.

상기 '135 특허는 $P$ 값이 사전 설정된 제 1 기준 값 및 제2 기준 값에 따라, $0 \le P < \text{제1 기준 값}$ 및 $\text{제2 기준 값} < P \le 1$ 구간에서는 제1 소리 신호를, 제1 기준 값 이상 제2 기준 값 이하 구간에서는 제2 소리 신호를 출력하는 출력 제어 로직을 제안한다.

한편, 상기 '135 특허의 딥러닝 알고리즘은 소음 혼합 신호로부터 음성 성분을 분리·추출하는 것과 이를 바탕으로한 판별적(discriminative) 접근에 해당하며, 입력 신호 내에 분리 가능한 음성 성분이 존재함을 전제로 한다. 
또한, 센서 포화에 의해 입력 파형의 충실도가 현저히 저하된 조건에서의 대응 수단은 별도로 개시하지 않는다.

\subsubsection*{2. 골전도 진동자를 이용한 조종사용 헬멧 (KR 20-2011-0004763 U)}

대한민국 공개실용신안 제20-2011-0004763호(이하 '763 실용신안)는 항공기 및 전차 조종사용 헬멧에 골전도 입출력 장치를 적용한 구조를 개시한다.
구체적으로, 상기 '763 실용신안은 사용자의 귀둘레 뼈에 밀착되는 귀마개(110)와 그 내부에 설치되는 골전도스피커(130), 골전도스피커와 귀마개 사이에 설치되어 골전도스피커가 귀둘레 뼈에 밀착 고정되도록 하는 탄성부재(150)를 포함하는 헬멧본체(100)를 개시한다.
또한, 사용자의 후두에 위치하여 음성에 의한 후두 진동을 전기적 신호로 변환하는 골전도마이크(170)와, 이를 후두 위치에 고정하기 위한 Y자 형상의 휨지지대(175)를 개시한다.

상기 '763 실용신안은 골전도 입출력 장치를 통해 귀마개로 외부 소음을 차단하면서도 골전도 경로로 음성 송수신이 가능하도록 하여, 조종사의 청력 보호와 교신 기능을 양립시키고자 하는 구성을 제안한다.

그러나 상기 '763 실용신안은 골전도 입출력 장치의 기계적 배치 및 밀착 고정 구조에 한정된 것으로서, 수신된 음성 신호에 대한 신호처리나 음성 향상 알고리즘을 개시하지 않는다. 또한 골전도 신호에 포함될 수 있는 신체 진동 잡음 성분의 처리나, 골전도 경로 고유의 제한된 주파수 응답을 보완하기 위한 대역 확장 수단에 대해서도 구체적 구성을 결여하고 있다.



\subsubsection*{3. 골진동 센서 및 마이크로폰 신호를 융합한 딥 러닝 음성 추출 및 노이즈 저감 방법 (KR 10-2429152 B1)}

대한민국 등록특허 제10-2429152호(이하 '152 특허)는 골진동 센서와 마이크로폰으로부터 각각 음성 신호를 샘플링하여 획득하고(Step 1), 골진동 센서 오디오 신호를 하이패스 필터링 모듈에 입력하여 직류 오프셋을 보정하고 저주파 노이즈를 제거한 후(Step 2), 이를 마이크로폰 오디오 신호와 심층 신경망(DNN) 모듈에 함께 입력하여 융합(fusing) 및 예측을 수행함으로써 노이즈 저감 후의 음성을 획득하는 방법을 개시한다(S3, S4).

상기 '152 특허는 골진동 신호의 저주파 대역 특성을 보완하기 위해, 하이패스 필터링을 거친 골진동 센서 오디오 신호에 대해 고주파 복원을 수행하여 대역 확장의 방법으로 주파수 범위를 확장한 후 심층 신경망 모듈에 입력하는 단계와, 대역 확장된 골진동 신호를 마이크로폰 신호 없이 직접 출력으로 사용하는 선택적 구성을 개시한다.
또한 심층 신경망 모듈의 구현 예로서 컨볼루션 순환 신경망 구조를 제시하고, 골진동 센서 오디오 신호 및 마이크로폰 오디오 신호를 각각 단기 푸리에 변환(STFT)하여 얻은 진폭 스펙트럼을 중첩하여 입력으로 사용하고, 예측 진폭 스펙트럼(EMS)을 출력하며, 타겟 진폭 스펙트럼(TMS)과 EMS에 대해 평균제곱오차(MSE)를 학습 목적함수로 사용하는 방식을 설명한다.

그러나 상기 '152 특허의 심층 신경망은 입력 스펙트럼으로부터 예측 스펙트럼(EMS)을 산출하는 추정방식에 해당하며, 입력 신호 내에 유효한 음성 성분이 존재함을 전제로 한다. 
센서 포화나 클리핑에 의해 입력 파형의 충실도가 현저히 저하된 조건에서의 음성 향상 수단은 별도로 다루지 않는다.


\subsubsection*{4. 헬멧을 위한 능동 소음 제거 시스템 (KR 10-2706589 B1)}

대한민국 등록특허 제10-2706589호(이하 '589 특허)는 헬멧의 정의된 공간 영역에서 특정 주파수 범위의 음압을 우선적으로 감쇠시키는 능동 소음 제거(ANC) 시스템을 개시한다. 
상기 '589 특허는 설치판, 라우드스피커, 참조 마이크로폰 및 ANC 제어 시스템을 포함하며, 참조 마이크로폰 신호에 기초하여 소음원과 역위상(antiphase)인 음파를 방출함으로써 조용한 구역(quiet zone)을 형성하는 피드-포워드 방식의 시스템을 개시한다.

상기 '589 특허의 배경기술은 바람 소음, 엔진 소음 등이 존재하는 모터사이클, 스카이다이빙 등 역동적인 스포츠 환경을 예시하며, 적응형 필터(LMS 알고리즘 등)를 사용할 수 있음을 개시한다.

그러나 상기 '589 특허는 참조 마이크로폰 신호에 기초한 역위상 음파 생성을 통해 소음을 상쇄하는 감쇠(cancellation) 방식으로서, 왜곡되거나 손실된 음성 신호를 복원하거나 음성 품질을 향상시키기 위한 별도의 음성 복원 기능은 개시하지 않는다. 
또한 제어 시스템으로 LMS 등 전통적 적응형 필터 방식을 예시할 뿐, 딥러닝 기반의 음성 처리 접근은 포함하지 않는다.



\subsubsection*{5. 골전도 통신이 가능한 안전모 (KR 20-0317070 Y1)}

대한민국 등록실용신안 제20-0317070호(이하 '070 실용신안)는 공사현장 등 소음이 심한 환경에서 사용하기 위한 골전도 통신 기능을 갖춘 안전모를 개시한다.
상기 '070 실용신안은 안전모의 내주면에 고정되어 이마에 접촉하는 골전도 마이크와, 관자놀이에 접촉하는 골전도 스피커, 및 이들을 연결하는 케이블로 이루어진 어셈블리를 구비하여, 귀마개를 착용하고 지내야 하는 산업 현장에서도 귀마개를 제거하지 않고 통신이 가능하도록 한다.

그러나 상기 '070 실용신안은 골전도 마이크 및 스피커의 물리적 배치 구조에 한정된 것으로서, 골전도 신호에 포함될 수 있는 신체 진동 잡음 성분의 제거나, 골전도 경로의 대역폭 한계를 보완하기 위한 신호처리 수단에 대해서는 구체적 구성을 결여하고 있다.


\subsubsection*{6. 골전도 센서를 구비한 청각 장치 (US 12340789 B2)}

미국 등록특허 제12,340,789호(이하 '789 특허)는 골전도 센서로부터 수집한 신호를 기반으로 합성 음성을 생성하는 청각 장치를 개시한다.
상기 '789 특허는 골전도 센서 신호를 음성 모델(speech model)의 제어 입력(control input)으로 사용하여 합성 음성 신호를 출력하는 구성을 개시하며, 상기 음성 모델로서 순환 신경망(RNN), WaveRNN 등의 자기회귀 음성 모델(autoregressive speech model)을 예시한다.
또한, 학습 단계에서 저소음 환경의 골전도 신호와 마이크로폰 신호 쌍을 이용하여 모델을 학습시키는 방법을 개시한다.

그러나 상기 '789 특허의 음성 모델은 골전도 센서 신호를 제어 입력으로 받아 새로운 음성 파형을 합성(synthesis)하는 구조로서, 골전도 신호가 유효한 음성 특징을 보존하고 있음을 암묵적으로 전제한다. 
특히, 극한 음압 환경에서 골전도 센서 자체에 포화나 클리핑이 발생하여 제어 입력의 신뢰성이 훼손된 경우, 이를 검출하고 손상된 입력으로부터 원본 음성을 복원하는 수단은 개시하지 않는다.



\subsection{비특허 문헌}

\begin{itemize}
    \item 기술적 배경 보강 목적
    \item 발명의 구체적 방법 및 추가적 아이디어가 확정되면, 필요한 내용을 조사 예정.
\end{itemize}

\newpage


\section{발명의 내용}

\subsection{해결하려는 과제}

본 발명은 상술한 종래 기술의 문제점 및 한계를 해소하기 위하여 창안된 것으로서, 극한 소음환경에서 사용자의 청력을 보호함과 동시에 명확한 음성 통신을 가능하게 하는 골전도 음향장치 및 인공지능 기반 음성 신호처리 기술을 헬멧과 일체화하여 제공하는 데 그 과제가 있다.

보다 구체적으로, 본 발명이 해결하고자 하는 과제는 다음과 같다.

\begin{enumerate}[label=(\alph*)]
    \item \textbf{[통합형 음성 통신 시스템 제공]}\\
    음성 통신 기능을 보호용 헬멧에 통합 적용함으로써, 실제 전장, 전투훈련, 공사 현장 등 극한 소음환경에서 사용자 상호 간의 명확한 음성 통신을 제공하는 동시에, 사용자의 청력을 보호하고 전반적인 안전성을 향상시키는 통합형 시스템을 제공하는 데 있다.

    \item \textbf{[청력 보호와 원활한 음성 통신]}\\
    청력 보호장비의 착용을 통해 고소음으로부터 청각 기관을 보호하면서, 귀를 통한 일반적인 공기전도 경로 대신 골전도 방식을 이용하여 사용자가 인지 가능한 음성 정보를 각 사용자에게 제공함으로써, 청력 보호와 원활한 통신을 동시에 달성하는 데 있다.
    
    \item \textbf{[선형 소음 제거 기술의 한계 극복]}\\
    비선형 왜곡 및 센서 포화가 발생하는 극심한 소음환경에서도 음성 신호의 명료도를 효과적으로 복원하여, 실시간 대화가 가능한 수준의 음성 통신 품질을 확보하는 데 있다. 이때, 인공지능 기반 음성 향상 모델과, 극한 소음 환경의 물리적 특성(비선형 왜곡, 센서 포화)을 반영한 학습 데이터 및 포화 마스크를 결합한 데이터 기반 접근에 의해, 선형·가우시안 잡음 모델에 기초한 종래 선형 필터링 기반 소음 제거 기술의 한계를 극복하는 것을 과제로 한다.
    % [수정] 지적 #8 대응: "물리적 특성을 반영한 신호처리 기법" → "데이터 기반 접근" 명시

    \item \textbf{[작업 수행 및 효율성 증대]}\\
    복합적이고 다양한 음원 신호(폭발음, 기계 소음, 통신 음성 등)가 혼재된 환경에서 소음 성분을 효과적으로 배제하고, 작업 및 전투 수행에 필요한 음성 정보만을 청명하게 추출·전달함으로써, 작전 수행 및 작업 간 효율성과 상황 인지 능력을 향상시키는 데 있다.

    \item \textbf{[경량·저지연 온디바이스(on-device) 인공지능 시스템 구축]}\\
    명료한 음성 신호를 제공하기 위한 신호처리 및 인공지능 기반 음성 추출·분리 프로세스를 헬멧 및 온디바이스(on-device) 단에서 구현함으로써, 외부 서버나 원격 처리를 필요로 하지 않고 현장에서 즉각적인 음향 알고리즘 수행과 인공지능 연산을 가능하게 하는 경량·저지연 시스템을 제공하는 데 있다.
\end{enumerate}

본 발명의 기술적 과제는 상기에서 언급한 내용에 한정되지 아니하며, 이하의 설명과 청구범위로부터 본 발명이 속하는 기술분야에서 통상의 지식을 가진 자라면 자명하게 도출할 수 있는 기타 과제들 역시 본 발명의 범위에 포함된다.


\subsection{과제의 해결 수단}

% ──────────────────────────────────────────────
% [신규 추가] 과제-해결수단 대응표
% ──────────────────────────────────────────────

상기 과제를 해결하기 위한 본 발명의 각 구성과 과제 간의 대응 관계는 하기 표~\ref{tab:task_solution_map}과 같다.

\begin{table}[!h]
    \centering
    \renewcommand{\arraystretch}{1.3}
    \begin{tabular}{c|l|l}
        \hline
        \textbf{과제} & \textbf{과제 내용} & \textbf{대응 구성} \\
        \hline\hline
        (a) & 통합형 음성 통신 시스템 제공 & S (시스템 총괄), A, B, C, D \\
        \hline
        (b) & 청력 보호와 원활한 음성 통신 & C (골전도 출력 및 청력 보호부), B \\
        \hline
        (c) & 선형 소음 제거 기술의 한계 극복 & A (멀티모달 음성 입력부), B (신호처리·AI 연산부) \\
        \hline
        (d) & 작업 수행 및 효율성 증대 & A, B, D (통신 인터페이스부) \\
        \hline
        (e) & 경량·저지연 온디바이스 AI 시스템 & B (신호처리·AI 연산부) \\
        \hline
    \end{tabular}
    \caption{과제와 해결 수단 구성의 대응 관계}
    \label{tab:task_solution_map}
\end{table}


\subsubsection{[S] 시스템 총괄}


\begin{enumerate}[label = \textbf{(S.\arabic*)}]

    \item 본 발명의 일 실시예에 따른 보호용 헬멧 일체형 음성 통신 장치는,
    골전도 센서 및 지향성 마이크 어레이를 포함하는 멀티모달 음성 입력부(A);
    입력 환경에 따라 연산 복잡도를 동적으로 조절하며 온디바이스 음성 향상을 수행하는 탈부착 가능한 신호처리·AI 연산부(B);
    사용자의 이도를 폐쇄하는 차음 구조체 및 골전도 진동자를 통해 음성과 방향 정보를 출력하는 골전도 음성 출력 및 청력 보호부(C);
    외부 통신 장치와 연동되는 통신 인터페이스부(D)를 포함한다.

    \item 상기 골전도 음성 출력 및 청력 보호부(C)는 상기 멀티모달
    음성 입력부(A)와 함께 헬멧 내피에 통합 배치되며, 상기 헬멧과의
    탈부착을 위한 결합 요소를 포함하는 것을 특징으로 한다.

    \item 상기 장치는 음압 레벨 130\,dB\,SPL 이상의 극한 소음 환경에서의 동작을 목적으로 구성되며, 현장 음성은 상기 멀티모달 음성 입력부(A)로부터 상기 신호처리·AI 연산부(B)를 거쳐 상기 골전도 음성 출력 및 청력 보호부(C)로 전달되고, 외부 통신 음성은 상기 통신 인터페이스부(D)로부터 상기 신호처리·AI 연산부(B)를 거쳐 상기 골전도 음성 출력 및 청력 보호부(C)로 전달되는, 두 경로의 음성 신호가 상기 신호처리·AI 연산부(B)에서 통합 처리되도록 구성된다. 상기 장치는, 상기 멀티모달 음성 입력부(A)의 공기전도 마이크로폰이 허용 입력 음압(AOP, 일 실시예에서 135\,dB\,SPL)을 초과하는 환경에서 센서 포화가 발생하는 것을 전제로 하며, 이에 대응하기 위한 포화 마스크 기반 적응적 융합(B-2) 및 골전도 우세 복원 모드(B-2)를 구비한다. 상기 골전도 센서는 두개골 전달 경로의 자연 감쇄에 의해 상기 공기전도 마이크로폰 대비 포화 발생 빈도가 현저히 낮으며, 이에 의해 공기전도 채널이 포화된 조건에서도 골전도 채널을 통한 음성 정보의 취득이 가능하다.
    % [신규 추가] 지적 #13 대응: 포화 전제 설계 논거, AOP 135dB, 골전도 감쇄 우위
\end{enumerate}


\begin{table}[!h]
    \centering
    \begin{tabular}{c|l|l}
        \hline
        구성 & 명칭 & 주요 역할 \\
        \hline\hline
        A & 멀티모달 음성 입력부
          & 골전도 센서, 지향성 마이크 어레이, 포화 마스크 생성 \\
        B & 신호처리·AI 연산부
          & 탈부착형 모듈, 적응형 연산 제어, 딥러닝 음성 향상 \\
        C & 골전도 음성 출력 및 청력 보호부
          & 차음 구조체, 골전도 진동자, 방향 정보 출력, 출력 리미터 \\
        D & 통신 인터페이스부
          & 군용 무전기 연동, 분대 내 근거리 통신, PTT \\
        \hline
    \end{tabular}
    \caption{발명의 주요 구성 요소}
    \label{tab:components}
\end{table}



\subsubsection{[A] 멀티모달 음성 입력부}

\paragraph{[A-1] 센서 및 신호 취득}

\begin{enumerate}[label = \textbf{(A.1.\arabic*)}]

    \item 본 발명의 멀티모달 음성 입력부는 사용자의 골조직 진동을 수집하는 골전도 센서 및 외부 음압을 수집하는 복수의 공기전도 마이크로폰을 포함하며, 상기 각 센서로부터 수집된 아날로그 신호를 개별적으로 처리하는 전처리 회로를 포함한다.
    % [수정] "공기전도마이크로폰" → "공기전도 마이크로폰" 용어 통일

    \item 상기 골전도 센서는 측두골, 전두골 또는 두정골 중 적어도 하나의 접촉 지점에 배치되며, 헬멧 내면 또는 착용 구조물에 의해 피부와의 접촉 상태를 유지하도록 구성된다.

    \item 상기 공기전도 마이크로폰은 사용자 구화 방향으로 지향성을 형성하는 빔포밍(Beamforming) 마이크 어레이로 구성되어 하드웨어 단계에서 신호 대 잡음비(SNR)를 일차적으로 개선한다. 
    또한, 상기 공기전도 마이크로폰은 귀폐쇄형 착용 구조에서 사용자에게 외부 환경 소리를 전달하기 위한 상황 인식(Situational Awareness) 경로의 입력원으로 기능할 수 있으며, 상기 경로의 음성 신호는 상기 신호처리·AI 연산부(B)를 거쳐 상기 골전도 음성 출력 및 청력 보호부(C)를 통해 사용자에게 전달된다.

\end{enumerate}


\paragraph{[A-2] 전처리 및 포화 마스크 생성}

\begin{enumerate}[label = \textbf{(A.2.\arabic*)}]

    \item 상기 전처리 회로는 각 채널별로 독립적인 대역 통과 필터 (Band Pass Filter, BPF)를 포함하여, 음성 대역 이외의 저주파 진동 및 고주파 전자기 노이즈를 제거한다. 
    특히, 상기 골전도 센서 채널은 신체 활동에 의한 저주파 노이즈를 차단하도록 설계되며, 상기 공기전도 마이크로폰 채널은 디지털 변환 시 앨리어싱(Aliasing) 방지를 위한 저역 통과(Low Pass) 특성을 포함한다.

    \item 상기 전처리된 각 채널의 신호는 동일한 클록 소스(Clock Source)에 의해 동기화된 멀티채널 ADC(Analog-to-Digital Converter)를 통해 디지털 데이터로 변환된다. 
    이를 통해 복수의 센서 간 위상(Phase) 정보가 보존되며, 보존된 위상 정보는 후단 신호처리·AI 연산부(B)에서의 채널 간 상관관계 분석 및 채널 정렬 처리에 활용된다.

    \item 상기 멀티모달 음성 입력부는 공기전도 및 골전도 신호의 진폭을 실시간 모니터링하여 포화 여부를 판단하는 포화 검출부를 포함한다. 
    상기 포화 검출부는 음압에 의한 공기전도 신호의 클리핑(Clipping) 및 물리적 충격에 의한 골전도 신호의 과도 응답 구간을 검출하여 시간-주파수 영역의 포화 마스크(Saturation Mask)를 생성한다. 상기 포화 검출은, 각 채널의 디지털 신호가 ADC 풀스케일의 기 설정된 비율(일 실시예에서 95\%)을 초과하는 샘플이 연속하여 기 설정된 개수(일 실시예에서 3개) 이상 검출되는 경우, 해당 시간 구간을 포화 구간으로 판정하는 것을 포함한다. 상기 시간-주파수 영역의 해상도는 상기 신호처리·AI 연산부(B)에서 사용하는 STFT 파라미터(일 실시예에서, 윈도우 길이 25\,ms, 홉 길이 10\,ms, NFFT 512)에 대응한다.
    % [수정] 지적 #4 대응: 포화 검출 판정 기준(ADC 95%, 연속 3샘플) 및 STFT 해상도 기재

    \item 상기 포화 검출부는, 골전도 센서와 공기전도 마이크로폰이 동시에 포화 상태인 수음 공백 구간을 검출하여, 해당 구간을 상기 신호처리·AI 연산부(B)에 통지하는 기능을 더 포함한다.

    \item 상기 포화 마스크는 시간-주파수 영역에서 신호의 왜곡 정도를 나타내며, 이진(Binary) 마스크 및 연속적(Soft) 마스크의 두 유형으로 생성된다. 상기 이진 마스크는, 시간 영역에서 상기 포화 판정 기준을 충족하는 구간을 검출하고, 해당 시간 구간에 대응하는 STFT 프레임의 전체 주파수 빈을 포화로 마킹하여 0 또는 1의 값을 부여한 것으로서, 상기 적응형 연산 제어 로직(B-1)의 모드 전환 판정에 사용된다. 상기 연속적 마스크는, 각 STFT 프레임 구간 내에서 시간 영역 신호가 상기 풀스케일 비율 임계치를 초과하는 샘플의 비율(클리핑 비율)을 산출하고, 상기 클리핑 비율을 시그모이드 함수에 의해 0$\sim$1의 연속값으로 변환하여, 해당 프레임의 전체 주파수 빈에 동일하게 적용한 것으로서, 상기 음성 향상 모델(B-2)의 적응적 융합 입력에 사용된다.
    % [수정] 지적 #4 대응: 이진/소프트 마스크 생성 로직·용도 분리 기재, 시간 영역 기반 검출

\end{enumerate}



\subsubsection{[B] 신호처리·AI 연산부}

\paragraph{[B-1] 적응형 연산 제어}

\begin{enumerate}[label = \textbf{(B.1.\arabic*)}]

    \item 본 발명의 신호처리·AI 연산부는 상기 멀티모달 음성 입력부(A)로부터 수신한 다채널 디지털 신호 및 상기 통신 인터페이스부(D)로부터 수신한 외부 통신 음성 신호를 처리하여 소음을 제거하고, 음성을 향상시키며, 입력 환경에 따라 연산 복잡도를 동적으로 조절하는 적응형 연산 제어 로직을 포함한다.

    \item 상기 적응형 연산 제어 로직은, 상기 멀티모달 음성 입력부(A)로부터 전달된 포화 마스크 정보 및 입력 신호의 외부 소음도(Noise Level)를 파라미터로 수신하여, 아래와 같이 다단계(Multi-stage)로 연산 모드를 전환하는 것을 특징으로 한다.

    \begin{itemize}
        \item \textbf{제0 모드(DSP 전용 모드):} 실시간 처리 지연이 기 설정된 지연 임계치(일 실시예에서 60\,ms)를
        초과하거나 프로세서 점유율이 기 설정된 부하 임계치(일 실시예에서 90\%) 이상인 경우, 딥러닝 모델을
        비활성화하고 고전적 신호처리(DSP) 알고리즘만을 이용하여 음성을
        처리한다.
        % [수정] 지적 #3 대응: 제0 모드 지연 60ms, CPU 90% 정량화

        \item \textbf{제1 모드(경량 AI 모드):} 포화 마스크가 비활성 상태(활성 비율 5\% 미만)이고
        외부 소음도가 기 설정된 제1 임계치 미만인 경우, 가용 연산 자원에 따라 경량 신경망(상기 신호처리·AI 연산부의 프로세서 점유율 50\% 이상 시) 또는
        고전적 신호처리(DSP) 알고리즘(50\% 미만 시)을 선택적으로 적용하여 전력 소모와 연산 부하를
        최소화하며 음성을 처리한다. 상기 제1 임계치는, 사용되는 공기전도 마이크로폰의 허용 입력 음압(AOP) 대비 소정의 마진을 확보하는 수준으로 설정되며, 일 실시예에서 AOP 135\,dB\,SPL 마이크 사용 시 약 115\,dB\,SPL로 설정된다.
        % [수정] 지적 #3 대응: 제1 모드 포화율 5%, 115dB SPL(AOP-20dB 마진), 경량NN/DSP 선택 기준

        \item \textbf{제2 모드(고성능 AI 모드):} 포화 마스크가 활성 상태(활성 비율 10\% 이상)이거나 외부 소음도가 상기 제1 임계치 이상인 경우, 고성능 딥러닝 기반 음성 향상 모델을 활성화하여 연산 자원을 집중 할당한다. 이때, 포화 마스크의 활성 비율이 소정 기준을 초과하는 경우, 상기 음성 향상 모델은 골전도 우세 복원 모드로 전환된다.
        % [수정] 지적 #3 대응: 제2 모드 포화율 10% 정량화
    \end{itemize}

    상기 모드 전환에는 채터링 방지를 위한 히스테리시스가 적용되며, 일 실시예에서 제2 모드에서 제1 모드로의 복귀 시 외부 소음도 기준을 110\,dB\,SPL 미만, 포화 마스크 활성 비율을 5\% 미만으로 설정한다. 상기 각 임계값은 운용 환경에 따라 조정 가능한 설계 파라미터이며, 구체적인 상태 전이 조건 및 전환 절차는 후술하는 실시예(§7.3.2)에서 상세히 기술한다.
    % [수정] 지적 #3 대응: 히스테리시스 조건(110dB/5% 복귀), 설계 파라미터 명시

    \item 상기 적응형 연산 제어 로직은, 입력 조건의 변화를 실시간으로
    모니터링하여 연산 모드 간 전환을 동적으로 수행하며, 모드 전환 시
    음성 처리의 연속성이 유지되도록 구성된다.

\end{enumerate}


\paragraph{[B-2] 음성 향상 모델 구조}

\begin{enumerate}[label = \textbf{(B.2.\arabic*)}]

    \item 상기 고성능 딥러닝 기반 음성 향상 모델은, 골전도 신호 및
    공기전도 신호의 시간 영역 또는 시간-주파수 영역 특징 표현과 상기
    포화 마스크를 입력으로 수신하여, 소음 및 왜곡이 억제된 음성 신호를
    출력한다.

    \item 상기 음성 향상 모델은, 입력 특징을 저차원 잠복 표현으로
    압축하는 인코더부, 상기 잠복 표현으로부터 잡음이 억제된 음성 특징을
    복원하는 디코더부, 및 인코더부와 디코더부 사이에서 다중 스케일
    정보를 보존하는 스킵 연결을 포함하는 인코더-디코더 구조로 구성된다. 일 실시예에서, 상기 인코더-디코더 구조는 복소수 영역(complex domain)에서 연산을 수행하며, 출력으로서 복소 비율 마스크(Complex Ratio Mask, CRM)를 추정할 수 있다. 상기 인코더-디코더 구조의 구체적 레이어 구성, 채널 수, 파라미터 규모 등은 후술하는 실시예(§7.3.3)에서 상세히 기술한다.
    % [수정] 지적 #5 대응: CRM 출력 방식을 "일 실시예"로 추가

    \item 상기 음성 향상 모델은 적응적 융합(Adaptive Fusion) 메커니즘을
    포함하여, 포화 마스크가 활성인 시간-주파수 구간에서는 골전도 신호
    특징의 가중치를 상대적으로 높이고, 비활성 구간에서는 공기전도 신호
    특징을 우세하게 활용함으로써, 센서별 신뢰도에 기반한 동적 채널
    융합을 수행한다.

    \item 상기 적응적 융합 메커니즘은, 포화 마스크가 
    기 설정된 포화율 임계치(일 실시예에서, 공기전도 채널의 시간-주파수 빈 중 활성 비율 70\% 이상) 이상의 시간-주파수 구간에서 활성 상태인 
    경우, 공기전도 신호 특징의 기여를 실질적으로 
    배제하고 골전도 신호 특징에 기반하여 음성을 
    복원하는 골전도 우세 복원 모드를 포함한다. 상기 포화율 임계치는 운용 환경에 따라 50$\sim$90\% 범위에서 조정 가능한 설계 파라미터이며, 판정 단위는 STFT 프레임 윈도우(일 실시예에서, 최근 10프레임) 내의 활성 비율로 한다.
    % [수정] 지적 #2 대응: "소정 비율" → 포화율 임계치 70% 정량화, 가변 범위 명시

    \item 상기 골전도 우세 복원 모드에서, 상기 
    디코더부는 골전도 신호의 유효 주파수 대역 상한을 
    초과하는 고주파 성분을 추정·보완하는 
    대역 확장(Bandwidth Extension)을 수행하여, 골전도 
    경로만으로도 음성 명료도를 유지한다. 
    상기 대역 확장은, 골전도 신호의 유효 주파수 대역 상한(일 실시예에서 약 4\,kHz)을 음성 명료도 확보에 필요한 대역(일 실시예에서 약 8\,kHz)까지 확장하는 것을 목표로 한다. 일 실시예에서, 상기 대역 확장은 4\,kHz 이하의 스펙트럼 포락(spectral envelope)으로부터 4$\sim$8\,kHz 대역의 스펙트럼 포락을 추정하는 경량 신경망(서브밴드 스펙트럼 포락 예측 방식)으로 구현되며, 상기 경량 신경망의 파라미터 규모는 일 실시예에서 약 300K이며, 목표 대역 및 모델 구조에 따라 가변된다. 상기 경량 신경망은 상기 음성 향상 모델의 디코더부에 부가되거나 독립 서브모듈로 구성될 수 있다. 구체적인 레이어 구성 및 훈련 방법은 후술하는 실시예(§7.3.3)에서 상세히 기술한다.
    % [수정] 지적 #2 대응: BWE 목표 대역(4→8kHz), 서브밴드 예측 방식, ~300K 파라미터 기재


    \item 상기 인코더부는 골전도 채널과 공기전도 채널의 특징을 각각
    독립적으로 인코딩한 후, 상기 적응적 융합 메커니즘에 의해 결합하여
    디코더부로 전달함으로써, 채널 간 상호보완적 정보를 최대한 보존하도록
    구성된다. 상기 듀얼 인코더 및 포화 마스크 기반 적응적 융합 구조는, 골전도 신호와 공기전도 신호를 단순 결합(concatenation)하여 단일 신경망에 입력하는 종래 구성과 구별된다. 종래 구성에서는 공기전도 채널이 포화된 경우에도 왜곡된 스펙트럼이 정상 신호와 동등하게 처리되는 반면, 본 발명의 구성에서는 포화 마스크가 왜곡 구간의 위치 및 신뢰도를 모델에 명시적으로 전달하여, 모델이 포화 구간에서 공기전도 정보를 의도적으로 배제하고 골전도 정보에 의존하는 학습이 가능하다.
    % [신규 추가] 지적 #5 대응: 선행기술('152 특허) 대비 구조적 차별점

\end{enumerate}


\paragraph{[B-3] 학습 데이터 및 견고성}

\begin{enumerate}[label = \textbf{(B.3.\arabic*)}]

    \item 상기 음성 향상 모델은, 포화, 클리핑, 및 비선형 왜곡 조건을 시뮬레이션에 의해 합성한 학습 데이터로 훈련되어, 실제 운용 환경에서의 신호 복원 견고성(Robustness)을 확보한다. 상기 합성 학습 데이터는 다음의 과정에 의해 생성된다: (i) 깨끗한 음성 데이터에 대해 골전도 전달 함수(저역 통과 특성, 일 실시예에서 차단 주파수 약 4\,kHz)를 적용하여 골전도 채널 신호를 시뮬레이션하고, (ii) 소음 데이터베이스로부터 선택된 소음을 소정의 SNR 범위(일 실시예에서 $-$10$\sim$+10\,dB)로 혼합하며, (iii) 센서 포화를 시뮬레이션하기 위해 하드 클리핑 함수($y = \mathrm{clip}(x, -T, +T)$, $T$: ADC 풀스케일) 또는 소프트 클리핑 함수($y = T \cdot \tanh(x/T)$)를 적용하고, (iv) 센서 비선형 응답을 시뮬레이션하기 위해 다항식 왜곡 모델($y = x + \beta_2 x^2 + \beta_3 x^3$)을 적용한다. 상기 비선형 계수 $\beta_2$, $\beta_3$는 사용 센서의 데이터시트 또는 캘리브레이션 측정에 의해 결정된다.
    % [수정] 지적 #6 대응: 합성 데이터 파이프라인 4단계 + 수식 기재

    \item 상기 학습 데이터는 상기 멀티모달 음성 입력부(A)의 구성에 대응하는 골전도 채널 및 공기전도 채널의 멀티채널 입력을 포함하며, 구체적인 훈련 방법 및 데이터 증강 기법은 후술하는 실시예(§7.3.5)에서 상세히 기술한다. 다만, 극한 음압 환경(120\,dB\,SPL 이상)에서의 실측 쌍 데이터 취득은 현실적으로 곤란하므로, 상기 학습 데이터는 주로 시뮬레이션 기반 합성에 의존하며, 합성 데이터와 실 운용 환경 간의 도메인 차이를 축소하기 위해 제한적으로 취득 가능한 저SPL 환경의 실측 데이터 또는 도메인 적응 기법이 추가 적용될 수 있다.
    % [수정] 지적 #6 대응: 도메인 갭 한계 솔직 기재 + 보완 수단

    \item 상기 배경기술(§3)에서 제시한 Westervelt 방정식 및 비선형 전파 특성은, 극한 소음 환경에서 선형 파동 방정식에 기초한 종래 필터링 기법이 적용 불가한 이유를 물리적으로 설명하기 위한 것이며, 본 발명의 해결 수단이 상기 비선형 방정식을 직접 풀거나 역산하는 것을 의미하지 않는다. 본 발명은, 비선형 전파 및 센서 포화가 신호에 미치는 영향을 학습 데이터에 시뮬레이션으로 반영하고, 포화 마스크가 왜곡 구간의 위치 및 신뢰도 정보를 모델에 명시적으로 전달함으로써, 딥러닝 모델이 비선형 왜곡 패턴에 대해 암묵적으로 학습하여 대응하도록 하는 데이터 기반 접근을 채택한다.
    % [신규 추가] 지적 #8 대응: 배경기술(Westervelt)과 해결수단(데이터 기반)의 논리 연결

\end{enumerate}



\paragraph{[B-4] 모델 경량화 및 실시간 처리}

\begin{enumerate}[label = \textbf{(B.4.\arabic*)}]

    \item 상기 음성 향상 모델은, 온디바이스 실시간 추론이 가능하도록
    모델 압축 및 경량화 기법이 적용되어, 제한된 연산 자원을 갖는
    임베디드 환경에서도 동작 가능하도록 구성된다.

    \item 상기 신호처리·AI 연산부는, 연산 복잡도와 음성 향상 성능 간의
    균형을 입력 조건에 따라 동적으로 조절하는 구조를 가지며, 전처리부터
    후처리에 이르는 전체 파이프라인이 실시간 양방향 통신 요건을 충족하도록
    구성된다.

\end{enumerate}


\paragraph{[B-5] 연산 하드웨어 및 신뢰성}

\begin{enumerate}[label = \textbf{(B.5.\arabic*)}]

    \item 상기 신호처리·AI 연산부는, 신경망 연산에 최적화된 프로세서를 포함하는 탈부착형 연산 모듈로 구현될 수 있으며, 사용 환경 또는 요구 성능에 따라 연산 모듈의 물리적 교체 및 모델 갱신이 가능하도록 구성된다.

    \item 상기 탈부착형 연산 모듈은, 전력 공급 및 데이터 인터페이스를 통해 상기 멀티모달 음성 입력부(A), 골전도 음성 출력 및 청력 보호부(C), 및 통신 인터페이스부(D)와 연결되며, 모듈 교체 시에도 시스템 인터페이스 호환성이 유지되도록 구성된다.

    \item 본 발명의 시스템은, 딥러닝 모델의 추론 실패 또는 연산 오류 발생 시, 각 센서 채널의 포화 상태에 따라 선택적으로 적용되는 폴백(Fallback) 경로를 구비한다. 상기 폴백 경로는, (i) 공기전도 채널이 정상인 경우에는 빔포밍 처리된 공기전도 신호를 우회 출력하고, (ii) 공기전도 채널이 포화 상태이고 골전도 채널이 정상인 경우에는 골전도 신호에 기본적인 전처리(대역 통과 필터링 및 이득 제어)를 적용하여 직접 출력하며, (iii) 양 채널이 동시에 포화 상태인 경우에는 무음 처리 또는 경고 신호를 출력하고 포화 해소를 대기하도록 구성된다.
    % [수정] 지적 #9 대응: 폴백 경로를 센서 포화 상태별 3분기로 재설계

    \item 상기 폴백 경로는, 상기 적응형 연산 제어 로직(B-1)과 연동하여 오류 감지 즉시 제0 모드로 복귀하도록 구성되며, 복귀 후 연산 환경이 정상화되면 이전 모드로 자동 전환된다. 또한, 상기 (iii)의 양 채널 동시 포화 상태에서는 통신 불능 상태를 사용자에게 시각적 또는 촉각적 수단으로 통지하는 구성을 더 포함할 수 있다.
    % [수정] 지적 #9 대응: 양 채널 동시 포화 시 사용자 통지 수단 추가
\end{enumerate}

\subsubsection{[C] 골전도 음성 출력 및 청력 보호부}

\paragraph{[C-1] 청력 보호 및 차음 구조}

\begin{enumerate}[label = \textbf{(C.1.\arabic*)}]

    \item 본 발명의 골전도 음성 출력 및 청력 보호부는 사용자의 이도를 폐쇄하여 외부 공기전도 소음으로부터 청력을 보호하는 차음 구조체를 포함하며, 상기 차음 구조체는 헬멧 내피에 일체로 형성되거나 탈착 가능하게 결합되도록 구성된다.

    \item 상기 차음 구조체에 의해 사용자의 이도가 폐쇄된 상태에서, 후술할 골전도 진동자는 두개골 진동을 통해 음성 신호를 와우에 직접 전달함으로써, 차음 상태를 유지하면서도 명료한 음성 수신을 가능하게 한다. 이때, 이도 폐쇄에 의한 폐쇄 효과(Occlusion Effect)는 주로 250\,Hz 이하의 저주파 대역에서 골전도 신호의 외이도 내 음압을 증가시키는 현상으로서, 골전도 출력의 저주파 감도를 향상시키는 한편, 음성 명료도에 기여하는 중고주파 대역(1$\sim$4\,kHz)에서는 그 효과가 제한적이다. 이에 따라, 상기 신호처리·AI 연산부(B)는 폐쇄 효과에 의한 저주파 과잉 증폭을 보정하기 위한 이퀄라이저를 골전도 출력 경로에 적용할 수 있다.
    % [수정] 지적 #10 대응: 폐쇄 효과 "활용" → "인지 및 보정" 전환, 주파수 한계 명시

    \item 상기 차음 구조체는 착용 환경 또는 임무 요건에 따라 차음 성능이 상이한 복수의 구조체로 교체 가능하도록 구성되어, 상황에 따른 차음 수준의 선택적 조절을 가능하게 한다.
\end{enumerate}


\paragraph{[C-2] 골전도 음성 출력 및 방향 정보}

\begin{enumerate}[label = \textbf{(C.2.\arabic*)}]

    \item 상기 골전도 음성 출력 및 청력 보호부는, 상기 신호처리·AI 연산부(B)로부터 출력된 향상된 음성 신호를 골전도 진동으로 변환하여 사용자의 두개골에 전달하는 좌측 및 우측 골전도 진동자를 포함하며, 상기 좌·우 진동자는 독립 채널로 구동된다.

    \item 상기 좌·우 골전도 진동자의 독립 채널 구동에 의해, 양이간 레벨차(ILD, Interaural Level Difference)에 기반한 좌우 방향 구분 정보를 포함하는 스테레오 골전도 출력이 제공될 수 있다. 다만, 골전도 경로에서는 두개골을 경유한 좌우 간 크로스토크(cross-talk)로 인해 공기전도 방식 대비 유효 ILD 범위가 제한되므로, 정밀한 음원 방위 판별보다는 대략적 좌우 방향 구분 수준의 정보 제공을 목적으로 한다.
    % [수정] 지적 #7 대응: "방향 인지" → "좌우 방향 구분" 하향, 크로스토크 한계 명시

    \item 상기 골전도 진동자는 골전도 전달 경로에서 발생하는 고주파 대역 감쇄를 보상하기 위해, 상기 신호처리·AI 연산부(B)에서 전달 함수 보정이 적용된 신호를 입력으로
    수신하도록 구성된다.

\end{enumerate}


\paragraph{[C-3] 출력 제한 및 청력 안전}

\begin{enumerate}[label = \textbf{(C.3.\arabic*)}]

    \item 상기 골전도 음성 출력 및 청력 보호부는, 골전도
    진동자의 출력이 기 설정된 안전 임계치를 초과하지 않도록
    실시간으로 제한하는 출력 리미터를 포함한다.

    \item 상기 출력 리미터는 충격성 소음에 대응하기 위해
    즉각적인 응답 특성을 갖도록 구성되며, 진동자 구동 신호의
    실시간 모니터링에 기초하여 출력을 제한함으로써 극한
    소음 환경에서의 골전도 경로를 통한 청력 손상을 방지한다.

    \item 상기 출력 리미터는 상기 신호처리·AI 연산부(B)의
    적응형 연산 제어 로직과 연동되어, 입력 소음 수준에 따라
    출력 제한 임계치가 동적으로 조절되도록 구성될 수 있다.

\end{enumerate}

\subsubsection{[D] 통신 인터페이스부}

\begin{enumerate}[label = \textbf{(D.\arabic*)}]

    \item 본 발명의 통신 인터페이스부는, 상기 신호처리·AI 연산부(B)로부터 출력된 향상된 음성 신호를 외부 통신 장치로 송신하고, 
    외부로부터 수신된 음성 신호를 상기 신호처리·AI 연산부(B)를 경유하여 상기 골전도 음성 출력 및 청력 보호부(C)로 전달하는 양방향 통신 경로를 제공한다. 
    이때, 송신되는 음성 신호는 상기 신호처리·AI 연산부(B)에 의해 소음 및 왜곡이 제거된 상태로 전송되므로, 수신 측에서는 별도의 소음 제거 처리 없이 명료한 음성을 수신할 수 있다.

\end{enumerate}


\paragraph{[D-1] 음성 송수신 인터페이스}

\begin{enumerate}[label = \textbf{(D.1.\arabic*)}]

    \item 상기 통신 인터페이스부는, 상기 향상된 음성 신호를
    군용 무전기 또는 외부 통신 단말에 전달하는 유선
    인터페이스를 포함하며, 군용 표준 오디오 인터페이스와
    호환되도록 구성된다.

    \item 상기 유선 인터페이스는 음성 신호를 소정의 음성 코덱으로 인코딩하여 송신하고, 외부로부터 수신된 인코딩 신호를 복호화하여 상기 신호처리·AI 연산부(B)로 전달하는 송수신 경로를 포함한다.

    \item 상기 통신 인터페이스부는 송신 경로와 수신 경로를 독립적으로 제어할 수 있는 구조를 가지며, 송신 측 음성 처리 상태와 무관하게 수신 음성이 상기 신호처리·AI 연산부(B)를 경유하여 상기 골전도 음성 출력 및 청력 보호부(C)를 통해 지속적으로 재생되도록 구성된다.

\end{enumerate}


\paragraph{[D-2] 보조 통신 및 사용자 제어}

\begin{enumerate}[label = \textbf{(D.2.\arabic*)}]

    \item 상기 통신 인터페이스부는, 외부 무전기 없이도
    인접한 복수의 장치 간 직접 음성 통신이 가능한 근거리
    무선 통신 채널을 더 포함하며, 상기 근거리 무선 통신
    채널은 상기 유선 인터페이스와 독립적으로 동작하도록
    구성된다.

    \item 상기 근거리 무선 통신 채널은 저지연 음성 전송
    요건을 충족하는 무선 통신 프로토콜을 이용하여 구현된다.
    일 실시예에서, 상기 근거리 무선 통신 채널은 LE Audio(BLE 5.2 이상, LC3 코덱 기반 CIS(Connected Isochronous Stream))를 이용하여 구현되며, 다른 실시예에서 DECT 또는 전용 Sub-GHz 무선 채널을 이용하여 구현될 수 있다. 상기 근거리 무선 통신 채널은 유선 통신이 불가한 환경에서의 보조 통신
    수단으로 기능한다.
    % [수정] 지적 #17 대응: LE Audio(CIS) 주 실시예 특정, "이에 준하는" 삭제

    \item 상기 통신 인터페이스부는 사용자가 송신 상태를
    직접 제어할 수 있는 PTT(Push-To-Talk) 스위치를
    포함하며, 상기 PTT 스위치는 본체로부터 분리된
    위치에 배치 가능하도록 구성된다.

\end{enumerate}


\subsubsection{[M] 방법 발명}

\paragraph{[M-1] 독립항 대응 기본 단계}

\begin{enumerate}[label = \textbf{(M.1.\arabic*)}]

    \item 본 발명의 일 실시예에 따른 보호용 헬멧 일체형
    음성 통신 장치를 이용한 음성 처리 방법은, 하기의
    단계를 포함한다.

    \begin{enumerate}[label = \textbf{(\alph*)}]

        \item 골전도 센서 및 공기전도 마이크로폰 어레이로부터 각각 골전도 음성 신호 및 공기전도 음성 신호를 취득하는 단계;

        \item 상기 공기전도 음성 신호 및 골전도 음성 신호의 진폭을 실시간 모니터링하여 포화 구간을 검출하고, 시간-주파수 영역의 포화 마스크를 생성하는 단계;

        \item 상기 포화 마스크 정보 및 입력 신호의 외부 소음도를 기반으로 연산 모드를 결정하는 단계;

        \item 상기 결정된 연산 모드에 따라, 상기 골전도 음성 신호, 공기전도 음성 신호 및 포화 마스크를 활용하여 소음 및 왜곡이 억제된 향상된 음성 신호를 출력하는 음성 향상 처리 단계;

        \item 상기 향상된 음성 신호를 골전도 진동으로 변환하여 사용자의 두개골에 전달하는 단계;

        \item 상기 향상된 음성 신호를 외부 통신 장치로 송신하는 단계; 및

        \item 상기 통신 인터페이스부(D)를 통해 외부로부터 수신된 음성 신호를 상기 신호처리·AI 연산부를 경유하여, 골전도 전달 경로 보정 및 출력 레벨 제어를 적용한 후 골전도 진동으로 변환·출력하는 단계.
        % [수정] 수신 경로 처리 내용 명시

    \end{enumerate}

    \item 상기 (c) 단계에서의 연산 모드 결정은, 실시간 처리 지연이 임계치를 초과하거나 가용 연산 자원이 임계치 미만인 경우 DSP 알고리즘만을 활용하는 제0 모드를 선택하고, 포화 마스크가 비활성 상태이고 외부 소음도가 제1 임계치 미만인 경우 경량 신경망 또는 DSP를 혼용하는 제1 모드를 선택하며, 포화 마스크가 활성 상태이거나 외부 소음도가 제1 임계치 이상인 경우 고성능 딥러닝 모델을 활성화하는 제2 모드를 선택하는 것을 포함한다.
    % [수정] "임계치" → "제1 임계치" B.1.2와 용어 통일


    \item 상기 (d) 단계에서의 음성 향상 처리는, 골전도 신호 및 공기전도 신호의 시간 영역 또는 시간-주파수 영역 특징 표현을 인코더-디코더 구조의 딥러닝 모델에 입력하여, 포화 마스크가 나타내는 신뢰도 정보에 기반한 신호 복원을 수행하는 것을 포함한다.

\end{enumerate}


\paragraph{[M-2] 종속항 대응 세부 단계}

\begin{enumerate}[label = \textbf{(M.2.\arabic*)}]

    \item 상기 (d) 단계는, 포화 마스크가 활성인
    시간-주파수 구간에서는 골전도 신호 특징의 가중치를
    상대적으로 높이고, 비활성 구간에서는 공기전도 신호
    특징을 우세하게 활용하는 적응적 융합을 수행하는
    단계를 더 포함한다.

    \item 상기 음성 처리 방법은, 딥러닝 모델의 추론
    실패 또는 연산 오류 발생 시, 공기전도 채널 및 골전도 채널의 포화 상태에 따라 (i) 빔포밍 처리된 공기전도 신호의 우회 출력, (ii) 골전도 신호의 직접 출력, 또는 (iii) 무음 처리 중 하나를 선택적으로 적용하는 폴백(Fallback)
    단계를 더 포함한다.
    % [수정] 지적 #9 대응: 방법항 폴백 3분기

    \item 상기 (e) 단계는, 차음 구조체에 의해 사용자의
    이도가 폐쇄된 상태에서 발생하는 폐쇄 효과(Occlusion Effect)를 고려하여, 저주파 과잉 증폭에 대한 보정이 적용된 골전도
    진동을 출력하는 단계를 포함한다.
    % [수정] 지적 #10 대응: 방법항 폐쇄 효과 "활용" → "보정 적용"

    \item 상기 (e) 단계는, 좌측 및 우측 골전도 진동자를
    독립 채널로 구동하여 양이간 레벨차(ILD)에 기반한 대략적 좌우 방향 구분 정보를 포함하는 스테레오 골전도 출력을
    제공하는 단계를 더 포함한다.
    % [수정] 지적 #7 대응: 방법항 ILD 하향

    \item 상기 (f) 단계는, 향상된 음성 신호를 군용
    표준 오디오 인터페이스를 통해 유선으로 송신하는
    단계; 및 외부 무전기 없이 인접한 복수의 장치
    간 근거리 무선 통신 채널을 통해 직접 송수신하는
    단계를 포함한다.

    \item 상기 음성 처리 방법은, 포화 및 비선형 왜곡
    조건을 포함하는 극한 소음 환경의 학습 데이터로
    훈련된 딥러닝 모델을 이용하여, 실제 운용 환경에서의
    신호 복원 견고성을 확보하는 단계를 더 포함한다.

    \item 상기 골전도 진동 출력 단계는, 골전도 전달 경로의 고주파 감쇄를 보상하기 위한 전달 함수 보정이 적용된 신호를 출력하는 단계를 더 포함한다.

    \item 상기 골전도 진동 출력 단계는, 골전도 진동자의 출력이 기 설정된 안전 임계치를 초과하지 않도록 실시간으로 제한하되, 입력 소음 수준에 따라 출력 제한 임계치를 동적으로 조절하는 단계를 더 포함한다.

    \item 상기 (d) 단계는, 포화 마스크의 활성 비율이 
    소정 기준을 초과하는 경우, 공기전도 신호의 기여를 
    배제하고 골전도 신호에 기반하여 음성을 복원하되, 
    골전도 신호의 제한된 주파수 대역을 보상하는 대역 
    확장을 수행하는 단계를 더 포함한다.

    \item 상기 음성 처리 방법은, 골전도 센서와 
    공기전도 마이크로폰이 동시에 포화 상태인 수음 
    공백 구간이 검출된 경우, 직전 유효 프레임의 
    STFT 진폭 스펙트럼을 유지(hold)하고 위상을 선형 외삽하는 스펙트럼 보간을 수행하되, 수음 공백 구간이 기 설정된 제1 프레임 수(일 실시예에서 3프레임, 30\,ms)를 초과하는 경우에는 점진적 페이드아웃을 적용하고, 기 설정된 제2 프레임 수(일 실시예에서 10프레임, 100\,ms)를 초과하는 경우에는 무음 처리하는 
    단계를 더 포함한다.
    % [수정] 지적 #11 대응: 보간 알고리즘(스펙트럼 hold+위상 외삽), 3단계 전환 기준 정량화

\end{enumerate}

\section{발명의 효과}

본 발명에 따른 보호용 헬멧 일체형 음성 통신 장치는 상술한 바와 같은 
구성 및 작용을 통하여 다음과 같은 효과를 제공한다.

\begin{enumerate}[label = \textbf{(\arabic*)}]

    \item 본 발명에 따르면, 골전도 음성 출력 및 청력 보호부(C)의 차음 구조체에 의해 외부 공기전도 소음으로부터 사용자의 청력을 보호하면서, 골전도 경로를 통해 명료한 음성을 전달할 수 있으므로, 음압 레벨 130\,dB\,SPL 이상의 극한 소음 환경에서도 청력 보호와 원활한 음성 통신을 동시에 달성할 수 있다. 
    
    \item 이에 따라, 종래의 귀마개 또는 귀덮개 착용에 수반되는 음성 소통 단절 문제가 해소되어, 사용자는 청력 손상 없이 지휘 통신, 주변 상황 인지, 및 사용자 간 의사소통 등 필수적인 음성 정보를 수용할 수 있다.

    \item 또한, 상기 멀티모달 음성 입력부(A)를 통해 수집된 현장 음성 및 환경음이 상기 신호처리·AI 연산부(B)에 의해 소음이 억제된 상태로 골전도 경로를 통해 사용자에게 전달되므로, 차음 구조체에 의해 이도가 폐쇄된 상태에서도 동료의 육성, 경고음 등 현장 상황 인지에 필요한 음성 정보를 수용할 수 있다. 다만, 상기 상황 인식 경로는 신호처리 파이프라인을 경유하므로 자연 청취 대비 소정의 처리 지연이 수반되며, 골전도 출력 경로의 대역폭 제한 및 방향 정보의 정확도에 한계가 존재한다.
    % [수정] 지적 #16 대응: SA 경로 한계 병기
    

    \item 또한, 멀티모달 음성 입력부(A)의 포화 검출부가 생성하는 포화 마스크를 신뢰도 정보로 활용하여, 신호처리·AI 연산부(B)가 센서별 신뢰도에 기반한 동적 채널 융합을 수행하므로, 센서 포화 및 비선형 왜곡이 발생하는 조건에서도 골전도 신호와 공기전도 신호의 상보적 특성을 최대한 보존하여 음성 명료도를 확보할 수 있다.

    \item 또한, 포화 마스크의 활성 비율에 기반하여 골전도 우세 복원 모드로 자동 전환되므로, 공기전도 센서가 완전 포화되는 극한 조건에서도 골전도 경로만을 이용한 음성 복원이 가능하여, 종래 다채널 융합 기술에서 대응하지 못하였던 완전 포화 시나리오에 대한 통신 연속성을 확보할 수 있다.
    
    \item 또한, 신호처리·AI 연산부(B)의 적응형 연산 제어 로직에 의해, 입력 환경의 소음 수준 및 포화 상태에 따라 연산 복잡도가 동적으로 조절되므로, 제한된 전력 및 연산 자원을 갖는 온디바이스 환경에서도 실시간 음성 향상이 가능하며, 전력 효율과 음성 품질 간의 균형을 유지할 수 있다. 
    
    \item 나아가, 딥러닝 모델의 추론 실패 또는 연산 오류 발생 시에는 상기 적응형 연산 제어 로직에 의해 DSP 기반 처리 모드로 자동 전환되어, 통신의 연속성이 보장된다.

    \item 또한, 상기 신호처리·AI 연산부(B)에 의해 소음 및 왜곡이 제거된 향상된 음성 신호가 통신 인터페이스부(D)를 통해 외부로 송신되므로, 수신 측 장비의 변경 없이도 명료한 음성이 전달되어, 기존 통신 인프라와의 호환성을 유지하면서 통신 품질이 향상된다.
    
    \item 아울러, 상기 통신 인터페이스부(D)는 외부 통신 장치와의 연동 및 근거리 직접 통신을 모두 지원하여, 다양한 운용 상황에서 유기적인 통신 연결을 제공한다.

    \item 또한, 상기 신호처리·AI 연산부(B)가 탈부착형 모듈로 구현되어 프로세서 교체 및 모델 갱신이 가능하므로, 기술 발전에 따른 성능 향상 및 다양한 운용 환경에 대한 적응이 용이하다. 
    
    \item 나아가, 음성 입출력부 및 청력 보호부가 헬멧에 통합됨으로써 별도의 장비 착용이 불필요하여, 사용자의 기동성과 작업 효율성이 유지된다.

    \item 또한, 딥러닝 모델의 추론 실패 시에도, 각 센서 채널의 포화 상태에 따라 공기전도 또는 골전도 신호를 선택적으로 우회 출력하는 폴백 경로가 구비되어, 시스템 장애 상황에서도 가용한 채널을 통한 통신의 연속성이 유지된다.
    % [수정] 지적 #9 대응: 효과 기재 폴백 3분기 반영

\end{enumerate}

다만, 본 발명에 따른 효과는 이상에서 예시된 내용에 의해 제한되지 않으며, 본 발명이 속하는 기술분야에서 통상의 지식을 가진 자가 본 명세서로부터 추론 가능한 모든 효과를 포함하는 것으로 이해되어야 한다.





\section{도안의 간단한 설명}

\begin{itemize}
    \item \textbf{도 1}: 시스템 전체 블록도. 멀티모달 음성 입력부(A), 신호처리·AI 연산부(B), 골전도 음성 출력 및 청력 보호부(C), 통신 인터페이스부(D)의 연결 관계와 포화 마스크 신호 흐름을 나타낸다.
    \item \textbf{도 2}: 적응형 연산 제어 상태 전이도. 제0 모드(DSP 전용), 제1 모드(경량 AI), 제2 모드(고성능 AI) 및 골전도 우세 복원 모드 간의 전환 조건과 폴백 경로를 나타낸다.
    \item \textbf{도 3}: 음성 향상 모델 아키텍처. 듀얼 인코더(골전도/공기전도), 포화 마스크 기반 적응적 융합 블록, 디코더, 및 대역 확장 서브모듈의 구조를 나타낸다.
    \item \textbf{도 4}: 오디오 파이프라인 흐름도. 전처리, 포화 검출 및 마스크 생성, 빔포밍, 모드별 추론, 후처리의 순서와 각 단계별 세부 처리를 나타낸다.
    \item \textbf{도 5}: 포화 시나리오별 신호 경로. 경우 1(공기전도 부분 포화), 경우 2(공기전도 완전 포화), 경우 3(양 채널 동시 포화)에 따른 신호 처리 경로 변화를 나타낸다.
\end{itemize}

\section{발명을 실시하기 위한 구체적인 내용}

% §7 발명을 실시하기 위한 구체적인 내용 — 전체 itemize
% Step 2 (subsection별 항목 나열 + 태깅) 완료본

\subsection{7.1 실시예 개관}
% [수정] 지적 #12 대응: 태그 → 본문 작성

이하에서는 본 발명의 일 실시예에 따른 보호용 헬멧 일체형 음성 통신 장치의 구체적인 구성 및 동작을 설명한다. 도 1을 참조하면, 본 실시예의 장치는 멀티모달 음성 입력부(A), 신호처리·AI 연산부(B), 골전도 음성 출력 및 청력 보호부(C), 및 통신 인터페이스부(D)로 구성된다.

\paragraph{(a) 신호 경로}

본 장치에는 두 가지 주요 음성 신호 경로가 존재한다. 현장 음성 경로는 멀티모달 음성 입력부(A)에서 수집된 골전도 및 공기전도 신호가 신호처리·AI 연산부(B)를 거쳐 골전도 음성 출력 및 청력 보호부(C)를 통해 사용자에게 전달되는 경로(A→B→C)이다. 외부 통신 음성 경로는 통신 인터페이스부(D)를 통해 외부로부터 수신된 음성 신호가 신호처리·AI 연산부(B)를 거쳐 골전도 음성 출력 및 청력 보호부(C)로 전달되는 경로(D→B→C)이다. 양 경로의 음성 신호는 신호처리·AI 연산부(B)에서 통합 처리된다.

\paragraph{(b) 동작 모드 개요}

신호처리·AI 연산부(B)는 입력 환경에 따라 제0 모드(DSP 전용), 제1 모드(경량 AI), 제2 모드(고성능 AI)의 3단계로 동작하며, 포화 마스크 활성 비율, 외부 소음도, 프로세서 점유율에 기반하여 모드 간 전환이 동적으로 수행된다. 각 모드의 구체적 전환 조건은 §7.3.2에서 상세히 기술한다.

\paragraph{(c) 모듈 구조}

신호처리·AI 연산부(B)는 탈부착 가능한 연산 모듈로 구현될 수 있으며, 전력 공급 및 데이터 인터페이스를 통해 상기 A, C, D 구성부와 연결된다. 모듈의 물리적 인터페이스 구조 및 환경 내구성은 §7.6에서 상세히 기술한다.

\subsection{7.2 멀티모달 음성 입력부(A)}

\subsubsection{7.2.1 골전도 센서 배치 및 스펙}
% [수정] 지적 #13, #12 대응: 태그 → 본문 작성

\paragraph{(a) 센서 유형 및 응답 대역}

골전도 센서는 압전식 가속도계(piezoelectric accelerometer)로 구성되며, 일 실시예에서 유효 응답 대역은 약 50$\sim$4{,}000\,Hz이다. 상기 대역은 골전도 경로를 통해 전달되는 음성 신호의 주요 에너지 분포 대역에 대응한다.

\paragraph{(b) 배치 위치}

상기 골전도 센서는 측두골, 전두골, 또는 하악골 중 적어도 하나의 접촉 지점에 배치된다. 측두골 배치는 헬멧 내면과의 접촉이 용이하고 음성 신호 전달 효율이 비교적 높으며, 전두골 배치는 헬멧 착용 시 안정적인 접촉 압력 유지에 유리하다.

\paragraph{(c) 탄성 마운트 구조}

상기 골전도 센서는 탄성 마운트를 통해 헬멧 내면에 고정되며, 사용자 두부와의 접촉 압력을 소정 범위(일 실시예에서 2$\sim$5\,N)로 유지하도록 구성된다. 상기 탄성 마운트는 헬멧 외부로부터의 구조 진동이 센서에 직접 전달되는 것을 억제하는 진동 절연 기능을 겸한다.

\paragraph{(d) 센서 포화 특성}

상기 압전식 가속도계의 최대 허용 가속도는 일 실시예에서 약 50$\sim$100\,g이다. 140\,dB\,SPL 이상의 극심한 충격(폭발 등) 시 골전도 센서에도 일시적 포화가 발생할 수 있으나, 두개골 전달 경로의 자연 감쇄(약 30$\sim$40\,dB)에 의해 공기전도 마이크로폰 대비 포화 발생 빈도가 현저히 낮다. 이는 본 발명의 골전도 우세 복원 모드(§5.2 B.2.4)의 기술적 근거이다.

\paragraph{(e) 비음성 저주파 억제}

신체 활동(보행, 저작 등)에 의한 저주파 잡음을 억제하기 위해, 골전도 센서 채널에 고역 통과 필터(일 실시예에서 차단 주파수 약 80$\sim$100\,Hz)가 적용된다.

\subsubsection{7.2.2 마이크 어레이 구성 및 빔포밍}
% [수정] 지적 #13, #12 대응: 태그 → 본문 작성

\paragraph{(a) MEMS 마이크 스펙}

공기전도 마이크로폰은 고AOP(Acoustic Overload Point) MEMS 마이크로폰으로 구성되며, 일 실시예에서 채널 수는 2$\sim$4개, 응답 대역은 약 100\,Hz$\sim$10\,kHz, AOP는 약 135\,dB\,SPL이다. 상기 AOP는 Table~\ref{tab:extreme_noise_examples}에 기재된 중공업 환경(110$\sim$140\,dB) 일부에서는 정상 동작을 보장하나, 전장 환경(140$\sim$170\,dB\textsuperscript{+})에서는 센서 포화가 불가피하다. 이는 본 발명의 포화 마스크 기반 적응적 융합 및 골전도 우세 복원 모드의 존재 이유이다.

\paragraph{(b) 어레이 배치}

상기 MEMS 마이크는 헬멧 외면의 사용자 구화 방향(전방 하부)에 배치되며, 마이크 간 간격은 일 실시예에서 약 2$\sim$4\,cm로 설정된다. 상기 배치에 의해 사용자 음성 방향으로의 지향성을 형성하고, 측방 및 후방 소음을 억제하는 빔포밍이 가능하다.

\paragraph{(c) 빔포밍}

빔포밍 알고리즘의 상세 및 SNR 개선 효과는 §7.3.1(c)에서 기술한다.

\paragraph{(d) EMI 차폐 및 풍잡음 방지}

상기 MEMS 마이크는 EMI(전자기 간섭) 차폐를 위한 금속 쉴드 케이스 내에 배치되며, 마이크 개구부에는 풍잡음 방지를 위한 방풍 스크린(windscreen)이 설치된다.

\paragraph{(e) 상황 인식 경로}

상기 공기전도 마이크 어레이는 음성 통신 외에, 외부 환경음을 수집하여 신호처리·AI 연산부(B)를 거쳐 골전도 출력부(C)를 통해 사용자에게 전달하는 상황 인식(Situational Awareness) 경로의 입력원으로도 기능한다.

\subsection{7.3 신호처리·AI 연산부(B)}

\subsubsection{7.3.1 오디오 파이프라인}
% [수정] 지적 #4, #11, #12 대응: 태그 → 본문 작성

본 실시예에서의 오디오 파이프라인은 전처리, 포화 검출 및 마스크 생성, 빔포밍, 딥러닝 추론, 및 후처리의 순서로 구성된다.

\paragraph{(a) 전처리 및 ADC}

멀티모달 음성 입력부(A)로부터 수신된 골전도 및 공기전도 아날로그 신호는, 동일한 클록 소스에 의해 동기화된 멀티채널 ADC를 통해 디지털 신호로 변환된다. 일 실시예에서, 샘플링 레이트는 16\,kHz, ADC 해상도는 24비트로 설정된다. 각 채널의 디지털 신호는 STFT(Short-Time Fourier Transform)에 의해 시간-주파수 영역으로 변환되며, STFT 파라미터는 일 실시예에서 윈도우 길이 25\,ms(400샘플, Hann 윈도우), 홉 길이 10\,ms(160샘플), NFFT 512로 설정된다. 이에 의해 시간-주파수 영역의 해상도는 100 프레임/초 $\times$ 257 주파수 빈이 된다.

\paragraph{(b) 포화 검출 및 마스크 생성}

포화 검출은 시간 영역에서 수행된다. 각 채널의 디지털 신호가 ADC 풀스케일의 95\%(일 실시예)를 초과하는 샘플이 연속 3개(일 실시예) 이상 검출되면, 해당 시간 구간을 포화 구간으로 판정한다.

이진 마스크는 상기 포화 구간에 대응하는 STFT 프레임의 전체 주파수 빈을 1(포화)로, 그 외 구간을 0(정상)으로 마킹한다.

소프트 마스크는 각 STFT 프레임 구간 내에서의 클리핑 비율(풀스케일 초과 샘플 수 / 프레임 내 총 샘플 수)을 산출하고, 시그모이드 함수 $\sigma(\cdot)$에 의해 0$\sim$1의 연속값으로 변환하여 해당 프레임의 전체 주파수 빈에 동일하게 적용한다.

\paragraph{(c) 빔포밍}

공기전도 마이크 어레이에 대해 지연-합(Delay-and-Sum, DS) 빔포밍 또는 MVDR(Minimum Variance Distortionless Response) 빔포밍이 적용된다. 빔포밍에 의해 일 실시예에서 약 6$\sim$10\,dB의 하드웨어 단계 SNR 개선이 기대된다. 빔포밍 출력은 후단 딥러닝 모델의 공기전도 채널 입력으로 사용된다.

\paragraph{(d) 딥러닝 추론}

적응형 연산 제어 로직(§7.3.2)에 의해 결정된 모드에 따라, 제0 모드에서는 빔포밍 및 위너 필터만 적용되고, 제1 모드에서는 경량 신경망 또는 DSP가 적용되며, 제2 모드에서는 고성능 음성 향상 모델(§7.3.3)이 적용된다.

\paragraph{(e) 후처리}

딥러닝 모델의 출력인 복소 비율 마스크(CRM)를 입력 복소 스펙트로그램에 적용하여 향상된 복소 스펙트로그램을 산출하고, iSTFT(Inverse STFT) 및 오버랩-애드(overlap-add) 방법에 의해 시간 영역 파형으로 변환한다. 변환된 파형은 DAC를 거쳐 골전도 진동자 구동 신호 및 통신 인터페이스 송신 신호로 출력된다.

\paragraph{(f) 양 채널 동시 포화 시 수음 공백 구간 처리}

골전도 센서와 공기전도 마이크로폰이 동시에 포화 상태인 수음 공백 구간에서는, 직전 유효 프레임의 STFT 진폭 스펙트럼을 유지(hold)하고 위상을 선형 외삽하는 스펙트럼 보간을 수행한다. 수음 공백이 3프레임(30\,ms)을 초과하는 경우에는 점진적 페이드아웃을 적용하고, 10프레임(100\,ms)을 초과하는 경우에는 무음 처리한다. 30\,ms 이하의 결손은 인간 청각의 시간 마스킹 특성에 의해 인지가 어려우며, 100\,ms 이상의 결손은 명확한 단절로 인지되므로 무음 처리가 적절하다.

\subsubsection{7.3.2 적응형 연산 제어 (+ 폴백)}
% [수정] 지적 #3, #9, #12 대응: 태그 → 본문 작성

\paragraph{(a) 모드 전환 조건}

적응형 연산 제어 로직의 모드 전환 조건은 다음 표와 같다. 상기 각 임계값은 일 실시예에서의 설계값이며, 운용 환경에 따라 조정 가능하다.

\begin{table}[H]
    \centering
    \renewcommand{\arraystretch}{1.3}
    \begin{tabular}{c|l|l|l}
        \hline
        \textbf{파라미터} & \textbf{제0 모드 진입} & \textbf{제1→제2 전환} & \textbf{제2→제1 복귀} \\
        \hline\hline
        외부 소음도 (SPL) & — & $\geq$ 115\,dB & $<$ 110\,dB \\
        \hline
        포화 마스크 활성율 & — & $\geq$ 10\% & $<$ 5\% \\
        \hline
        프로세서 점유율 & $\geq$ 90\% & — & — \\
        \hline
        파이프라인 지연 & $>$ 60\,ms & — & — \\
        \hline
        모델 오류 & 감지 시 즉시 & — & — \\
        \hline
    \end{tabular}
    \caption{모드 전환 조건 (일 실시예)}
\end{table}

상기 프로세서 점유율은 상기 신호처리·AI 연산부(B) 모듈 단독의 프로세서 점유율을 기준으로 한다. 제1 임계치(115\,dB\,SPL)는 사용 마이크의 AOP(일 실시예에서 135\,dB\,SPL) 대비 20\,dB 마진으로서, 센서 포화 시작 전에 고성능 모드로의 선제 전환을 위한 설계값이다.

\paragraph{(b) 히스테리시스 조건}

모드 전환 시 채터링(반복 전환)을 방지하기 위해, 상향 전환(제1→제2)과 하향 전환(제2→제1)의 임계치에 히스테리시스를 적용한다. 일 실시예에서, 제1→제2 전환은 외부 소음도 115\,dB 이상 또는 포화율 10\% 이상에서 발생하고, 제2→제1 복귀는 외부 소음도 110\,dB 미만 및 포화율 5\% 미만이 동시에 충족될 때 수행된다. 제0→제1 복귀는 프로세서 점유율이 80\% 미만으로 감소하고 파이프라인 지연이 40\,ms 미만으로 회복된 상태가 기 설정된 시간(일 실시예에서 1초) 이상 지속된 경우 수행된다.

\paragraph{(c) 모드 전환 시 음성 연속성}

모드 전환 시 출력 음성의 불연속을 방지하기 위해, 전환 직전 모드의 출력과 전환 직후 모드의 출력에 대해 크로스페이드(cross-fade)를 적용한다. 일 실시예에서 크로스페이드 구간은 2프레임(20\,ms)으로 설정된다.

\paragraph{(d) 제1 모드 내 경량 신경망/DSP 선택}

제1 모드에서는 상기 신호처리·AI 연산부의 프로세서 점유율이 50\% 이상인 경우 경량 신경망(INT8 양자화 모델)을 적용하고, 50\% 미만인 경우 고전적 DSP 알고리즘(빔포밍 + 위너 필터)만을 적용한다.

\paragraph{(e) 폴백 동작}

딥러닝 모델의 추론 실패 또는 연산 오류가 감지된 경우, 적응형 연산 제어 로직은 즉시 제0 모드로 복귀한다. 오류 감지 조건은 (i) 모델 추론의 출력값이 기 설정된 범위를 초과하는 이상치 발생, (ii) 추론 지연이 기 설정된 한도(일 실시예에서 60\,ms)를 초과, 또는 (iii) 런타임 예외(exception) 발생을 포함한다.

제0 모드 복귀 후의 폴백 출력은 센서 포화 상태에 따라 결정된다:
\begin{itemize}
    \item 공기전도 채널 정상: 빔포밍 처리된 공기전도 신호를 우회 출력
    \item 공기전도 채널 포화 + 골전도 채널 정상: 골전도 신호에 대역 통과 필터링 및 이득 제어를 적용하여 직접 출력
    \item 양 채널 동시 포화: 무음 처리 또는 경고 신호 출력, 포화 해소 대기
\end{itemize}

정상 복귀 조건은 오류 원인 해소 후 프로세서 점유율 및 파이프라인 지연이 정상 범위 내로 회복된 것이 확인된 경우이며, 제0→제1 복귀 히스테리시스 조건을 충족한 후 이전 모드로 자동 전환된다.

\subsubsection{7.3.3 음성 향상 모델 아키텍처}
% [수정] 지적 #5, #12 대응: 태그 → 본문 작성

본 실시예에서의 음성 향상 모델은 인코더-디코더 구조를 기반으로 하며, 골전도 채널과 공기전도 채널을 각각 독립적으로 인코딩하는 듀얼 인코더 구성을 채택한다.

\paragraph{(a) 전체 구조}

상기 모델은 골전도 인코더, 공기전도 인코더, 적응적 융합 블록(§7.3.4에서 후술), 디코더, 및 대역 확장 서브모듈로 구성된다. 인코더와 디코더 사이에는 다중 스케일 정보를 보존하기 위한 스킵 연결이 포함된다.

\paragraph{(b) 듀얼 인코더}

골전도 인코더 및 공기전도 인코더는 각각 5층의 복소 컨볼루션(Complex Conv2D) 블록으로 구성되며, 각 블록은 컨볼루션 레이어, 배치 정규화(Batch Normalization), 및 활성화 함수(PReLU)를 포함한다. 채널 수는 하위 레이어로부터 상위 레이어로 16 → 32 → 64 → 128 → 256으로 점진적으로 증가한다. 두 인코더는 구조가 동일하나 파라미터를 공유하지 않으며, 각 채널의 고유한 특징 표현을 독립적으로 학습한다.

\paragraph{(c) 병목(Bottleneck)}

인코더 출력은 2층의 복소 LSTM(Complex LSTM, hidden dimension = 256)으로 구성된 병목부를 통과하여 시간적 의존성을 모델링한다.

\paragraph{(d) 디코더}

디코더는 5층의 복소 전치 컨볼루션(Complex ConvTranspose2D) 블록으로 구성되며, 인코더와 대칭적 구조를 갖는다. 채널 수는 256 → 128 → 64 → 32 → 16으로 감소한다. 각 디코더 레이어에는 대응하는 인코더 레이어의 스킵 연결이 결합된다.

\paragraph{(e) 입출력 형태}

입력은 각 채널(골전도, 공기전도)의 STFT 복소 스펙트로그램 및 포화 마스크(소프트 마스크)로 구성된다. 일 실시예에서 샘플링 레이트 16\,kHz, NFFT 512 기준 각 채널의 입력 텐서 형태는 (배치, 2, 시간 프레임, 257)이며, 여기서 2는 실수부 및 허수부를 나타낸다. 포화 마스크는 (배치, 1, 시간 프레임, 257)의 형태로 융합 블록에 입력된다. 출력은 복소 비율 마스크(CRM)이며, 입력 공기전도 스펙트로그램에 CRM을 적용하여 향상된 복소 스펙트로그램을 산출한 후 iSTFT에 의해 시간 영역 파형으로 변환된다.

\paragraph{(f) 손실 함수}

학습 시 손실 함수는 SI-SNR(Scale-Invariant Signal-to-Noise Ratio) 손실과 STOI(Short-Time Objective Intelligibility) 손실의 가중 합으로 구성된다. 일 실시예에서 $\mathcal{L} = \mathcal{L}_{\text{SI-SNR}} + 0.3 \cdot \mathcal{L}_{\text{STOI}}$ 를 사용한다.

\paragraph{(g) 파라미터 규모}

상기 듀얼 인코더, 병목, 디코더, 및 적응적 융합 블록을 포함한 전체 모델의 파라미터 수는 일 실시예에서 약 2.1M이다.

\paragraph{(h) 골전도 단독 복원 및 대역 확장 서브모듈}

골전도 우세 복원 모드(§5.2 B.2.4)가 활성화된 경우, 공기전도 인코더의 출력은 적응적 융합 블록에서 실질적으로 배제되고, 디코더는 골전도 인코더의 출력만을 이용하여 음성을 복원한다. 이때, 골전도 신호의 유효 대역(약 4\,kHz)을 초과하는 고주파 성분을 보완하기 위해 대역 확장(BWE) 서브모듈이 동작한다.

상기 BWE 서브모듈은, 디코더 출력의 4\,kHz 이하 스펙트럼 포락으로부터 4$\sim$8\,kHz 대역의 스펙트럼 포락을 추정하는 경량 신경망으로 구성된다. 일 실시예에서, 상기 경량 신경망은 2층의 완전 연결(FC) 레이어(hidden dimension = 256, 활성화 함수 ReLU) 및 1층의 1D 컨볼루션 레이어(커널 크기 3, 출력 채널 수 = 목표 주파수 빈 수)로 구성되며, 입력으로 4\,kHz 이하 대역의 로그 멜 스펙트럼 포락을, 출력으로 4$\sim$8\,kHz 대역의 로그 스펙트럼 포락을 추정한다. 파라미터 규모는 일 실시예에서 약 300K이다.

\subsubsection{7.3.4 적응적 융합 메커니즘}
% [수정] 지적 #15, #12 대응: 태그 → 본문 작성

상기 적응적 융합 블록은, 골전도 인코더 출력과 공기전도 인코더 출력을 포화 마스크에 기반하여 동적으로 결합하는 모듈이다.

\paragraph{(a) 구현 방식: 어텐션 게이트}

일 실시예에서, 적응적 융합은 어텐션 게이트(Attention Gate) 구조로 구현된다. 포화 마스크(소프트 마스크, 0$\sim$1 연속값)가 어텐션 가중치의 산정에 입력되어, 각 시간-주파수 구간에서의 채널별 기여도를 결정한다. 구체적으로, 포화 마스크 값이 높은 구간(포화 심화)에서는 골전도 인코더 출력의 가중치가 상대적으로 증가하고, 낮은 구간(정상)에서는 공기전도 인코더 출력의 가중치가 우세해지도록 구성된다. 상기 어텐션 가중치는 포화 마스크와 각 인코더 출력의 특징을 입력으로 하는 1$\times$1 컨볼루션 및 시그모이드 활성화 함수에 의해 산출된다.

\paragraph{(b) 융합 위치}

상기 적응적 융합은 각 인코더의 최종 출력(병목 입력 직전)에서 수행되어, 융합된 특징이 병목 및 디코더로 전달된다. 추가적으로, 스킵 연결 경로에서도 각 스케일별로 동일한 포화 마스크 기반 가중치 조절이 적용되어 다중 스케일 융합이 이루어질 수 있다.

\paragraph{(c) 골전도 단독 경로 전환}

포화 마스크의 활성 비율이 기 설정된 포화율 임계치(일 실시예에서 70\%)를 초과하는 경우, 상기 어텐션 게이트는 공기전도 인코더 출력의 가중치를 실질적으로 0으로 설정하여 골전도 인코더 출력만이 디코더로 전달되는 골전도 단독 경로로 전환된다. 이 전환은 어텐션 가중치의 연속적 변화에 의해 이루어지므로, 모드 전환 시 출력 음성의 불연속이 방지된다.

\paragraph{(d) 선행기술 대비 구조적 차별점}

종래 구성(예: 선행기술 '152 특허의 CRNN)에서는 골전도 신호와 공기전도 신호의 스펙트럼을 단순 결합(concatenation)하여 단일 신경망에 입력하며, 포화 여부에 대한 정보가 모델에 전달되지 않는다. 이 경우, 공기전도 채널이 포화 상태에 있을 때 왜곡된 스펙트럼이 정상 신호와 동등하게 처리되어 모델 출력의 품질이 저하된다. 본 발명의 포화 마스크 기반 적응적 융합 구조에서는, 포화 마스크가 왜곡 구간의 위치 및 신뢰도를 모델에 명시적으로 전달하므로, 모델이 포화 구간에서 공기전도 정보를 의도적으로 배제하고 골전도 정보에 의존하는 처리가 가능하다. 이러한 왜곡 전파의 구조적 차단은 포화 시나리오를 고려하지 않는 종래 단순 융합 방식으로부터는 달성할 수 없는 효과이다.

\subsubsection{7.3.5 학습 데이터 생성 및 훈련 방법}
% [수정] 지적 #6, #12 대응: 태그 → 본문 작성

\paragraph{(a) 합성 데이터 생성 파이프라인}

본 발명의 음성 향상 모델은, 극한 소음 환경에서의 실측 쌍 데이터(깨끗한 음성과 극한 소음 동시 녹음) 취득이 현실적으로 곤란하므로, 시뮬레이션 기반 합성 데이터를 이용하여 훈련된다. 합성 데이터 생성 파이프라인은 다음의 단계로 구성된다.

\begin{enumerate}
    \item \textbf{깨끗한 음성 데이터 확보}: 공개 음성 데이터셋(예: LibriSpeech, VCTK, KsponSpeech 등)으로부터 다양한 화자 및 발화 조건의 깨끗한 음성을 확보한다.
    
    \item \textbf{골전도 채널 시뮬레이션}: 상기 깨끗한 음성에 골전도 전달 함수를 적용하여 골전도 채널 신호를 합성한다. 일 실시예에서 상기 골전도 전달 함수는 차단 주파수 약 4\,kHz의 저역 통과 특성으로 근사하되, 실측 골전도 전달 함수가 확보된 경우 이를 대체 적용하여 공진 피크 및 반공진(anti-resonance) 특성을 반영할 수 있다. 또한, 신체 활동에 의한 저주파 잡음(호흡, 심박 등)을 가산하여 현실성을 향상시킨다.
    
    \item \textbf{소음 혼합}: 소음 데이터베이스로부터 선택된 소음을 소정의 SNR 범위(일 실시예에서 $-$10$\sim$+10\,dB)로 공기전도 채널 신호에 혼합한다. 소음 데이터베이스는 군용 소음(총성, 폭발음, 차량 소음 등), 산업 소음(해머, 프레스, 발파 등), 및 광대역 환경 소음을 포함한다.
    
    \item \textbf{센서 포화 시뮬레이션}: 공기전도 채널에 대해 하드 클리핑($y = \mathrm{clip}(x, -T, +T)$, $T$: ADC 풀스케일) 또는 소프트 클리핑($y = T \cdot \tanh(x/T)$)을 적용한다. 클리핑 임계치 $T$는 데이터 증강 목적으로 $\pm$30\% 범위에서 무작위 변동시킨다.
    
    \item \textbf{센서 비선형 응답 시뮬레이션}: 다항식 왜곡 모델($y = x + \beta_2 x^2 + \beta_3 x^3$)을 적용하여 센서의 비선형 응답을 시뮬레이션한다. 비선형 계수 $\beta_2$, $\beta_3$는 사용 센서의 데이터시트 또는 캘리브레이션 측정에 의해 결정된다.
\end{enumerate}

\paragraph{(b) 멀티채널 데이터 구성}

학습 데이터의 각 샘플은 (i) 골전도 채널 시뮬레이션 신호, (ii) 소음 및 포화가 적용된 공기전도 채널 신호, (iii) 포화 마스크(§7.3.1의 방법에 의해 생성), 및 (iv) 목표(target) 깨끗한 음성 신호의 4-튜플로 구성된다. 골전도 및 공기전도 채널은 동일 시간축으로 동기화된다.

\paragraph{(c) 데이터 증강}

훈련 데이터의 다양성을 확보하기 위해, SNR 변조(소음 레벨 무작위 조절), 소음 유형 혼합(정상 소음 + 충격성 소음 중첩), 포화율 변조(클리핑 임계치 변동), 및 골전도 전달 함수의 파라미터 변동(차단 주파수 $\pm$500\,Hz 변동)을 적용한다.

\paragraph{(d) 합성 데이터의 한계}

상기 합성 데이터는 실 운용 환경과의 도메인 차이(domain gap)가 존재할 수 있으며, 특히 골전도 전달 함수의 단순 LPF 근사는 개인별 두개골 형상 및 센서 접촉 조건에 따른 변동을 충분히 반영하지 못할 수 있다. 이를 보완하기 위해, 제한적으로 취득 가능한 저SPL 환경의 실측 데이터를 이용한 파인튜닝 또는 도메인 적응 기법이 추가 적용될 수 있다.

\subsubsection{7.3.6 모델 경량화 (모드별 차등 전략)}
% [수정] 지적 #12 대응: 태그 → 본문 작성

\paragraph{(a) 모드별 양자화 전략}

상기 음성 향상 모델은 적응형 연산 제어의 모드에 따라 상이한 양자화 수준이 적용된다.

\begin{itemize}
    \item \textbf{제1 모드(경량 AI)}: INT8 양자화가 적용된 경량 모델을 사용한다. INT8 양자화에 의해 모델 크기가 약 1/4로 축소되며, 정수 연산에 의한 추론 속도 향상이 달성된다. 음성 향상 성능은 FP16 대비 소폭 저하될 수 있으나, 저소음 환경(제1 모드 동작 조건)에서는 충분한 품질이 유지된다.
    \item \textbf{제2 모드(고성능 AI)}: FP16(반정밀도 부동소수점) 모델을 사용한다. 포화 및 고소음 환경에서의 음성 향상 성능을 최대화하기 위해 모델의 전체 표현력을 활용한다.
\end{itemize}

\paragraph{(b) 추론 하드웨어 및 지연}

일 실시예에서, 상기 신호처리·AI 연산부의 프로세서는 NVIDIA Jetson Orin NX(FP16 기준 약 50\,TOPS) 또는 동등 성능의 임베디드 NPU를 사용한다. TensorRT 등의 추론 최적화 엔진을 적용하며, 상기 음성 향상 모델(약 2.1M 파라미터, FP16)의 단일 프레임 추론 시간은 일 실시예에서 약 5\,ms(추정)이다.

\paragraph{(c) 전체 파이프라인 지연 목표}

전처리(STFT 윈도우 + 홉: 약 25\,ms), 추론(약 5\,ms), 후처리(iSTFT + 오버랩-애드: 약 10\,ms)를 포함한 전체 파이프라인의 종단 간(end-to-end) 지연 목표는 일 실시예에서 약 40\,ms로 설정된다. 이는 실시간 양방향 음성 통신에서 사용자가 인지 가능한 지연 한계(약 100$\sim$150\,ms) 이내이다.

\subsection{7.4 골전도 음성 출력 및 청력 보호부(C)}

\subsubsection{7.4.1 차음 구조체}
% [수정] 지적 #12 대응: 태그 → 본문 작성

\paragraph{(a) 구조체 유형}

차음 구조체는 사용자의 이도를 폐쇄하여 외부 공기전도 소음의 유입을 차단하는 것으로서, 일 실시예에서 탄성 이어팁(in-ear tip), 커스텀 이어몰드(custom ear mold), 또는 헬멧 내장형 이어컵(ear cup) 중 하나 이상의 유형으로 구현된다.

\paragraph{(b) 차음 성능}

상기 차음 구조체의 차음 성능은 일 실시예에서 NRR(Noise Reduction Rating) 25\,dB 이상으로 설정된다. 차음 성능은 구조체 유형, 재질, 및 사용자 이도 형상에 따라 변동될 수 있다.

\paragraph{(c) 재질 및 밀폐 방식}

일 실시예에서 탄성 이어팁은 의료용 실리콘 또는 메모리 폼(memory foam)으로 제작되며, 이중 플랜지(double flange) 또는 삼중 플랜지 구조에 의해 이도와의 밀폐성을 향상시킨다.

\paragraph{(d) 교체 구조}

상기 차음 구조체는 착용 환경 또는 임무 요건에 따라 차음 성능이 상이한 복수의 구조체로 교체 가능하도록 구성된다. 일 실시예에서 구조체는 스냅핏(snap-fit) 방식에 의해 헬멧 내피의 소정 위치에 탈착된다.

\paragraph{(e) 헬멧 결합}

상기 차음 구조체와 헬멧 내피의 결합은 스냅핏, 나사, 또는 자석 결합 방식 중 하나로 구현되며, 현장에서의 신속한 교체가 가능하도록 공구 없이 탈착 가능한 구조가 바람직하다.

\subsubsection{7.4.2 골전도 진동자 스펙 및 구동 회로}
% [수정] 지적 #7, #12 대응: 태그 → 본문 작성

\paragraph{(a) 진동자 유형}

골전도 진동자는 선형 공진 액추에이터(LRA, Linear Resonant Actuator) 또는 압전 진동자로 구성된다. LRA는 음성 대역에서의 광대역 응답 특성 및 소비 전력 효율이 우수하며, 압전 진동자는 고주파 대역의 응답 특성이 상대적으로 우수하다.

\paragraph{(b) 응답 대역}

상기 골전도 진동자의 주 응답 대역은 일 실시예에서 약 100$\sim$4{,}000\,Hz이며, 4{,}000\,Hz 이상의 대역에서는 진동자의 기계적 특성에 의해 점진적 감쇄가 발생한다. 약 8{,}000\,Hz까지의 감쇄 대역에서도 소정의 출력이 가능하나 효율이 현저히 저하된다.

\paragraph{(c) 배치 및 접촉 구조}

좌·우 골전도 진동자는 각각 좌·우 측두골 위치에 배치되며, 탄성 패드를 통해 두개골과의 접촉 압력(일 실시예에서 약 3$\sim$6\,N)을 유지한다.

\paragraph{(d) 구동 회로}

구동 회로는 DAC → 파워 앰프 → 진동자의 경로로 구성되며, 좌·우 진동자는 독립 채널로 구동된다.

\paragraph{(e) ILD 기반 방향 정보 및 한계}

좌·우 독립 채널 구동에 의해 양이간 레벨차(ILD) 기반의 대략적 좌우 방향 구분 정보를 제공할 수 있다. 다만, 골전도 경로에서는 두개골을 경유한 좌우 간 크로스토크가 존재하며, 좌우 간 감쇄는 주파수 의존적으로 약 0$\sim$15\,dB 수준이다. 이에 의해 유효 ILD 범위는 공기전도 방식(고주파에서 최대 약 20\,dB) 대비 현저히 제한되며, 특히 저주파 대역($<$1\,kHz)에서는 크로스토크가 커서 유효한 ILD 형성이 곤란하다. 따라서, 상기 방향 정보는 정밀한 음원 방위 판별이 아닌 대략적 좌우 방향 구분 수준에 한정된다. 또한, 양이간 시간차(ITD)의 정밀한 재현은 골전도 전달 경로의 특성상 곤란하므로, ITD 기반 방향 인지는 본 실시예의 범위에 포함하지 않는다.

\subsubsection{7.4.3 골전도 전달 특성 보정 (전달 함수 + 폐쇄 효과)}
% [수정] 지적 #10, #12 대응: 태그 → 본문 작성

\paragraph{(a) 골전도 전달 함수 측정 및 보정}

골전도 진동자로부터 와우까지의 전달 경로는 두개골의 기계적 특성에 의해 주파수 의존적 감쇄를 나타내며, 특히 고주파 대역(2\,kHz 이상)에서 감쇄가 현저하다. 이를 보상하기 위해, 상기 신호처리·AI 연산부(B)에서 골전도 전달 함수의 역특성에 기반한 이퀄라이저(역필터)를 설계하여 출력 신호에 적용한다. 상기 전달 함수는 대표적 두개골 모델에 기초한 평균 전달 함수를 사용하거나, 개인별 캘리브레이션에 의해 보정될 수 있다.

\paragraph{(b) 폐쇄 효과(Occlusion Effect)의 영향 및 보정}

차음 구조체에 의해 이도가 폐쇄된 상태에서는, 폐쇄 효과에 의해 주로 250\,Hz 이하의 저주파 대역에서 골전도 신호의 외이도 내 음압이 약 10$\sim$20\,dB 증가한다. 이는 골전도 출력의 저주파 감도를 향상시키나, 음성 명료도에 기여하는 중고주파 대역(1$\sim$4\,kHz)에서는 유의미한 효과가 제한적이다. 오히려 저주파 과잉 증폭이 음성 명료도를 저해할 수 있으므로, 상기 이퀄라이저 설계 시 폐쇄 효과에 의한 저주파 증폭을 감안하여 저주파 대역의 이득을 하향 보정한다.

\paragraph{(c) 대역 확장 신호에 대한 보정}

골전도 우세 복원 모드에서 대역 확장 서브모듈(§7.3.3(h))에 의해 생성된 4$\sim$8\,kHz 대역의 신호는, 골전도 진동자의 고주파 감쇄 특성이 가장 현저한 대역에 해당한다. 이에 따라 상기 대역에 대해서는 이퀄라이저의 고주파 부스트 이득을 강화하되, 출력 리미터(§7.4.4)에 의한 안전 제한 범위 내에서 조절된다.

\subsubsection{7.4.4 출력 리미터}
% [수정] 지적 #12 대응: 태그 → 본문 작성

\paragraph{(a) 안전 기준}

골전도 진동자의 출력 안전 기준은 ISO 389-3을 참조하여, 골전도 출력을 Force Level(dB re 1\,$\mu$N) 또는 가속도 레벨(dB re 1\,m/s$^2$)로 규정한다. 기 설정된 안전 임계치를 초과하는 출력은 출력 리미터에 의해 즉시 제한된다.

\paragraph{(b) 리미터 구현}

일 실시예에서, 출력 리미터는 진동자 구동 전압을 실시간 모니터링하여 기 설정된 최대 구동 전압을 초과하지 않도록 제한하는 구동 전압 제한 방식으로 구현된다. 다른 실시예에서, 진동자에 인접 배치된 가속도 센서의 피드백에 기초하여 실제 출력 가속도를 모니터링하고 제한하는 방식으로 구현될 수 있다.

\paragraph{(c) 응답 특성}

충격성 소음에 대응하기 위해, 상기 출력 리미터의 공격 시간(attack time)은 일 실시예에서 1\,ms 이하, 릴리즈 시간(release time)은 약 50$\sim$100\,ms로 설정된다.

\paragraph{(d) 적응형 연산 제어와의 연동}

상기 출력 리미터의 안전 임계치는 상기 신호처리·AI 연산부(B)의 적응형 연산 제어 로직과 연동되어, 입력 소음 수준에 따라 동적으로 조절될 수 있다. 일 실시예에서, 외부 소음도가 높은 환경(제2 모드)에서는 안전 임계치를 소정 범위 내에서 상향 조절하여 음성 명료도를 확보하고, 저소음 환경(제1 모드)에서는 기본 안전 임계치를 유지한다.

\subsection{7.5 통신 인터페이스부(D)}

\subsubsection{7.5.1 군용 무전기 유선 연동}
% [수정] 지적 #12 대응: 태그 → 본문 작성

\paragraph{(a) 커넥터 규격}

상기 통신 인터페이스부의 유선 인터페이스는, 군용 표준 오디오 커넥터(일 실시예에서 U-283/U 6-pin 또는 동등 규격)를 사용하여 외부 무전기 또는 통신 단말과 연결된다.

\paragraph{(b) 음성 코덱}

송신 음성 신호는 소정의 음성 코덱으로 인코딩되어 전송된다. 일 실시예에서 군용 환경에서는 MELPe 2400\,bps(NATO STANAG 4591), 범용 환경에서는 Opus 코덱이 적용될 수 있다.

\paragraph{(c) 향상 음성의 코덱 전 투입 이점}

상기 신호처리·AI 연산부(B)에 의해 소음 및 왜곡이 제거된 향상된 음성 신호가 코덱 인코딩 전에 투입되므로, 코덱이 소음 성분에 비트를 낭비하지 않아 동일 비트레이트에서의 음성 품질이 향상되고, 수신 측에서 별도의 소음 제거 처리가 불필요하다.

\subsubsection{7.5.2 근거리 무선 통신}
% [수정] 지적 #17, #12 대응: 태그 → 본문 작성

\paragraph{(a) 프로토콜}

일 실시예에서, 근거리 무선 통신 채널은 LE Audio(BLE 5.2 이상, LC3 코덱 기반 CIS(Connected Isochronous Stream))를 이용하여 구현된다. 다른 실시예에서 DECT 또는 전용 Sub-GHz 무선 채널을 이용하여 구현될 수 있다.

\paragraph{(b) 통신 사양}

일 실시예에서 통신 범위는 약 50\,m(개활지 기준), 최대 동시 접속 노드 수는 약 12개, 음성 전송 지연은 약 20$\sim$30\,ms로 설정된다.

\paragraph{(c) 보조 통신 수단으로서의 위치}

상기 근거리 무선 통신 채널은 유선 통신(§7.5.1)이 불가한 환경에서의 보조 또는 비상 통신 수단으로 기능하며, 유선 인터페이스와 독립적으로 동작한다.

\subsubsection{7.5.3 PTT 및 사용자 제어}
% [수정] 지적 #12 대응: 태그 → 본문 작성

\paragraph{(a) PTT 스위치}

상기 통신 인터페이스부는 사용자가 송신 상태를 직접 제어할 수 있는 PTT(Push-To-Talk) 스위치를 포함한다. 상기 PTT 스위치는 헬멧 외면에 배치되거나, 유선 또는 무선에 의해 본체로부터 분리된 위치(예: 총기 그립, 조끼 등)에 배치 가능하다.

\paragraph{(b) 채널 선택}

일 실시예에서, 사용자는 유선 통신 채널과 근거리 무선 통신 채널 간의 전환을 물리적 스위치 또는 소프트웨어 인터페이스에 의해 제어할 수 있다.

\subsubsection{7.5.4 수신 음성 경로 (D→B→C)}
% [수정] 지적 #12 대응: 태그 → 본문 작성

\paragraph{(a) 수신 경로 처리 순서}

외부로부터 수신된 음성 신호는 통신 인터페이스부(D)에서 코덱 디코딩된 후, 신호처리·AI 연산부(B)를 경유하여 다음의 순서로 처리된다: (i) 코덱 아티팩트 보정(저비트레이트 코덱에 의한 양자화 잡음 억제), (ii) 골전도 전달 함수 보정(§7.4.3의 이퀄라이저 적용), (iii) 출력 리미터(§7.4.4)에 의한 안전 레벨 제한. 처리된 신호는 골전도 음성 출력 및 청력 보호부(C)를 통해 사용자에게 전달된다.

\paragraph{(b) 현장 음성과의 믹싱}

현장 음성 경로(A→B→C)의 출력과 수신 음성 경로(D→B→C)의 출력은 신호처리·AI 연산부(B)에서 믹싱되어 골전도 출력부(C)로 전달된다. 일 실시예에서, 수신 음성에 우선순위를 부여하여 수신 음성 재생 시 현장 음성의 레벨을 소정량 감쇄(ducking)하는 방식이 적용될 수 있다.

\paragraph{(c) 독립 경로 구조}

송신 경로의 음성 처리 상태(음성 향상 모델의 추론 상태, 모드 전환 등)와 무관하게, 수신 음성은 독립적인 처리 경로를 통해 지속적으로 재생되도록 구성된다. 이에 의해 송신 측 처리 장애 시에도 수신 음성의 청취가 중단되지 않는다.

\subsection{7.6 모듈 구조, 환경 내구성 및 전원 관리}
% [수정] 지적 #12 대응: 태그 → 본문 작성

\paragraph{(a) 탈부착 모듈 커넥터}

신호처리·AI 연산부(B)의 탈부착형 연산 모듈은, 일 실시예에서 MIL-DTL-38999 호환 방수 다핀 커넥터를 통해 헬멧 본체에 결합된다. 상기 커넥터는 전력 공급 및 디지털 데이터(멀티채널 오디오 입출력, 제어 신호)를 단일 인터페이스로 통합하여 전달한다.

\paragraph{(b) 모듈 교체 및 업그레이드}

상기 탈부착형 연산 모듈은 현장에서 공구 없이 교체 가능하도록 구성되며, 모듈 교체 시에도 커넥터 규격의 호환성이 유지되어 상이한 성능 등급의 프로세서 모듈 또는 갱신된 모델이 탑재된 모듈로의 교체가 가능하다.

\paragraph{(c) 환경 내구성}

본 장치는 군용 및 중공업 환경에서의 사용을 목적으로 하며, 다음의 환경 내구성 요건을 충족하도록 설계된다.

\begin{itemize}
    \item 방진·방수: IP67 이상
    \item 충격·진동 내구성: MIL-STD-810H 준거
    \item 동작 온도: $-40 \sim +60$°C
    \item EMI/EMC: MIL-STD-461G 참조
\end{itemize}

\paragraph{(d) 전원 설계}

일 실시예에서, 신호처리·AI 연산부(B)의 프로세서(Jetson Orin NX 또는 동등)의 소비전력은 약 15$\sim$25\,W이며, 센서·진동자·통신부를 포함한 시스템 전체 소비전력은 약 20$\sim$35\,W로 추정된다. 리튬이온 배터리(일 실시예에서 약 100\,Wh)를 사용하는 경우, 연속 동작 시간은 약 3$\sim$5시간으로 산출된다.

\paragraph{(e) 전원 보호 및 모드 연동}

전원 보호 회로는 과충전·과방전 보호 기능을 포함하며, 배터리 잔량이 기 설정된 저전압 임계치 이하로 감소한 경우 적응형 연산 제어 로직에 의해 제0 모드(DSP 전용)로 강제 전환되어 잔여 동작 시간을 연장한다.

\subsection{7.7 변형예}
% [수정] 지적 #12 대응: 태그 → 본문 작성

\subsubsection{7.7.1 다른 웨어러블 적용}

본 발명의 구성은 보호용 헬멧에 한정되지 않으며, 다른 웨어러블 형태로 변형 적용될 수 있다.

이어마운트형 변형예에서는, 헬멧 없이 귀걸이형 골전도 장치에 멀티모달 음성 입력부(A) 및 골전도 음성 출력 및 청력 보호부(C)가 탑재되고, 신호처리·AI 연산부(B)는 벨트 착용형 별도 모듈로 구현될 수 있다.

페이스실드 통합형 변형예에서는, 용접·연마 작업용 페이스실드에 상기 A, C 구성부가 통합 배치되어, 안면 보호와 음성 통신을 동시에 제공할 수 있다.

\subsubsection{7.7.2 다른 AI 모델 대체}

§7.3.3에서 기술한 인코더-디코더 구조는 일 실시예에 해당하며, Conv-TasNet, DCCRN, FullSubNet, TF-GridNet 등 다른 비생성형(non-generative) 음성 향상 모델로 대체될 수 있다. 이 경우에도 포화 마스크 기반 적응적 융합 메커니즘(§7.3.4) 및 골전도 우세 복원 모드의 적용은 동일하게 유지된다.

나아가, 향후 경량 생성형 모델(예: 1$\sim$2 step diffusion 모델)의 임베디드 실시간 처리가 가능해지는 경우, 이를 적용하여 음성 품질을 더욱 향상시킬 수 있는 가능성이 있다. 다만, 현 시점에서 생성형 모델의 임베디드 실시간 추론은 연산량 및 지연 측면에서 실현이 곤란하므로, 본 실시예에서는 판별적 모델을 주 실시예로 채택한다.

\subsubsection{7.7.3 민수 적용}

본 발명의 구성은 군용 환경에 한정되지 않으며, 민수 분야에도 적용 가능하다. 건설·중공업 현장에서는 산업용 안전모에 본 발명의 구성을 적용하고, 통신 인터페이스부(D)를 산업용 무전기 규격에 맞게 변형할 수 있다. 항공·우주 분야에서는 발사체 주변 작업자, 항공기 정비사 등 극심한 소음에 노출되는 작업자를 위한 음성 통신 장비로 적용될 수 있다.

\subsection{7.8 부호의 설명}
% [수정] 지적 #12 대응: 태그 → 본문 작성

본 명세서에서 사용된 주요 참조부호는 다음과 같다.

\begin{table}[H]
    \centering
    \renewcommand{\arraystretch}{1.2}
    \begin{tabular}{c|l|l}
        \hline
        \textbf{부호} & \textbf{명칭} & \textbf{대응 도면} \\
        \hline\hline
        A & 멀티모달 음성 입력부 & 도 1, 도 4 \\
        \hline
        B & 신호처리·AI 연산부 & 도 1, 도 2, 도 3, 도 4 \\
        \hline
        C & 골전도 음성 출력 및 청력 보호부 & 도 1 \\
        \hline
        D & 통신 인터페이스부 & 도 1 \\
        \hline
    \end{tabular}
    \caption{주요 구성부 참조부호 일람}
\end{table}

% 통계 섹션 삭제 — 본문 작성 완료에 따라 불필요

\subsection{실시예}

\subsection{부호의 설명}




\part{청구 범위}
% [신규 추가] 지적 #14 대응: 청구범위 전문 작성

\section*{청구항}

\begin{enumerate}[label=\textbf{[\arabic*]}]

% ═══════════════════════════════════════════════
% 독립항 1 — 장치
% ═══════════════════════════════════════════════

    \item 보호용 헬멧에 일체화되는 음성 통신 장치로서,

    사용자의 골조직 진동을 수집하는 골전도 센서 및 외부 음압을 수집하는 복수의 공기전도 마이크로폰을 포함하고, 상기 공기전도 마이크로폰 및 골전도 센서 신호의 진폭을 실시간 모니터링하여 포화 구간을 검출하고 시간-주파수 영역의 포화 마스크를 생성하는 멀티모달 음성 입력부(A);

    상기 멀티모달 음성 입력부(A)로부터 수신한 다채널 디지털 신호를 처리하여 소음 및 왜곡이 억제된 향상된 음성 신호를 출력하되, 상기 포화 마스크에 기반하여 골전도 신호와 공기전도 신호의 채널별 기여도를 동적으로 조절하는 적응적 융합 메커니즘을 포함하는 인코더-디코더 구조의 음성 향상 모델과, 상기 포화 마스크 정보 및 외부 소음도에 기반하여 연산 복잡도를 동적으로 조절하는 적응형 연산 제어 로직을 포함하는 신호처리·AI 연산부(B);

    사용자의 이도를 폐쇄하여 외부 소음으로부터 청력을 보호하는 차음 구조체, 및 상기 신호처리·AI 연산부(B)로부터 출력된 향상된 음성 신호를 골전도 진동으로 변환하여 사용자의 두개골에 전달하는 골전도 진동자를 포함하는 골전도 음성 출력 및 청력 보호부(C); 및

    상기 향상된 음성 신호를 외부 통신 장치로 송신하고, 외부로부터 수신된 음성 신호를 상기 신호처리·AI 연산부(B)를 경유하여 상기 골전도 음성 출력 및 청력 보호부(C)로 전달하는 통신 인터페이스부(D)

    를 포함하는 음성 통신 장치.

% ═══════════════════════════════════════════════
% 종속항 2~10 — 장치
% ═══════════════════════════════════════════════

    \item 제1항에 있어서,
    상기 적응적 융합 메커니즘은, 상기 포화 마스크가 기 설정된 포화율 임계치 이상의 시간-주파수 구간에서 활성 상태인 경우, 공기전도 신호의 기여를 실질적으로 배제하고 골전도 신호에 기반하여 음성을 복원하는 골전도 우세 복원 모드를 포함하며,
    상기 골전도 우세 복원 모드에서 상기 음성 향상 모델의 디코더부는 골전도 신호의 유효 주파수 대역 상한을 초과하는 고주파 성분을 추정·보완하는 대역 확장을 수행하는 것을 특징으로 하는 음성 통신 장치.

    \item 제1항에 있어서,
    상기 적응형 연산 제어 로직은,
    실시간 처리 지연이 기 설정된 지연 임계치를 초과하거나 프로세서 점유율이 기 설정된 부하 임계치 이상인 경우 고전적 신호처리 알고리즘만을 이용하는 제0 모드;
    상기 포화 마스크가 비활성 상태이고 외부 소음도가 기 설정된 제1 임계치 미만인 경우 경량 신경망 또는 고전적 신호처리 알고리즘을 선택적으로 적용하는 제1 모드; 및
    상기 포화 마스크가 활성 상태이거나 외부 소음도가 상기 제1 임계치 이상인 경우 고성능 딥러닝 모델을 활성화하는 제2 모드
    를 포함하여 다단계로 연산 모드를 전환하는 것을 특징으로 하는 음성 통신 장치.

    \item 제3항에 있어서,
    상기 모드 전환에는 채터링 방지를 위한 히스테리시스가 적용되어, 상향 전환과 하향 전환의 임계치가 상이하게 설정되는 것을 특징으로 하는 음성 통신 장치.

    \item 제1항에 있어서,
    상기 포화 마스크는 이진(Binary) 마스크 및 연속적(Soft) 마스크의 두 유형으로 생성되며,
    상기 이진 마스크는 시간 영역에서 포화 판정 기준을 충족하는 구간에 대응하는 STFT 프레임의 전체 주파수 빈을 포화로 마킹하여 상기 적응형 연산 제어 로직의 모드 전환 판정에 사용되고,
    상기 연속적 마스크는 각 STFT 프레임 구간 내의 클리핑 비율을 연속값으로 변환하여 상기 음성 향상 모델의 적응적 융합 입력에 사용되는 것을 특징으로 하는 음성 통신 장치.

    \item 제1항에 있어서,
    상기 신호처리·AI 연산부(B)는, 상기 음성 향상 모델의 추론 실패 또는 연산 오류 발생 시, 각 센서 채널의 포화 상태에 따라
    공기전도 채널이 정상인 경우에는 빔포밍 처리된 공기전도 신호를 우회 출력하고,
    공기전도 채널이 포화 상태이고 골전도 채널이 정상인 경우에는 골전도 신호를 직접 출력하며,
    양 채널이 동시에 포화 상태인 경우에는 무음 처리 또는 경고 신호를 출력하는 폴백 경로를 구비하는 것을 특징으로 하는 음성 통신 장치.

    \item 제1항에 있어서,
    상기 골전도 음성 출력 및 청력 보호부(C)는, 상기 골전도 진동자의 출력이 기 설정된 안전 임계치를 초과하지 않도록 실시간으로 제한하는 출력 리미터를 포함하며,
    상기 출력 리미터는 충격성 소음에 대응하기 위한 즉각적 응답 특성을 갖는 것을 특징으로 하는 음성 통신 장치.

    \item 제1항에 있어서,
    상기 골전도 음성 출력 및 청력 보호부(C)는, 좌측 및 우측 골전도 진동자가 독립 채널로 구동되어 양이간 레벨차(ILD)에 기반한 좌우 방향 구분 정보를 포함하는 스테레오 골전도 출력을 제공하는 것을 특징으로 하는 음성 통신 장치.

    \item 제1항에 있어서,
    상기 신호처리·AI 연산부(B)는, 신경망 연산에 최적화된 프로세서를 포함하는 탈부착형 연산 모듈로 구현되며, 사용 환경 또는 요구 성능에 따라 연산 모듈의 물리적 교체 및 모델 갱신이 가능하도록 구성되는 것을 특징으로 하는 음성 통신 장치.

    \item 제1항에 있어서,
    상기 통신 인터페이스부(D)는, 외부 무전기 없이 인접한 복수의 장치 간 직접 음성 통신이 가능한 근거리 무선 통신 채널을 더 포함하며,
    상기 근거리 무선 통신 채널은 저지연 무선 통신 프로토콜을 이용하여 구현되고, 상기 유선 인터페이스와 독립적으로 동작하는 것을 특징으로 하는 음성 통신 장치.

% ═══════════════════════════════════════════════
% 독립항 11 — 방법
% ═══════════════════════════════════════════════

    \item 보호용 헬멧에 일체화된 음성 통신 장치를 이용한 음성 처리 방법으로서,

    (a) 골전도 센서 및 공기전도 마이크로폰 어레이로부터 각각 골전도 음성 신호 및 공기전도 음성 신호를 취득하는 단계;

    (b) 상기 공기전도 음성 신호 및 골전도 음성 신호의 진폭을 실시간 모니터링하여 포화 구간을 검출하고, 시간-주파수 영역의 포화 마스크를 생성하는 단계;

    (c) 상기 포화 마스크 정보 및 입력 신호의 외부 소음도를 기반으로 연산 모드를 결정하는 단계;

    (d) 상기 결정된 연산 모드에 따라, 상기 골전도 음성 신호, 공기전도 음성 신호 및 포화 마스크를 인코더-디코더 구조의 음성 향상 모델에 입력하되, 상기 포화 마스크에 기반하여 골전도 신호와 공기전도 신호의 채널별 기여도를 동적으로 조절하는 적응적 융합을 수행하여, 소음 및 왜곡이 억제된 향상된 음성 신호를 출력하는 음성 향상 처리 단계;

    (e) 상기 향상된 음성 신호를 골전도 진동으로 변환하여, 차음 구조체에 의해 이도가 폐쇄된 상태의 사용자의 두개골에 전달하는 단계;

    (f) 상기 향상된 음성 신호를 외부 통신 장치로 송신하는 단계; 및

    (g) 외부로부터 수신된 음성 신호를 골전도 전달 경로 보정 및 출력 레벨 제어를 적용한 후 골전도 진동으로 변환·출력하는 단계

    를 포함하는 음성 처리 방법.

% ═══════════════════════════════════════════════
% 종속항 12~20 — 방법
% ═══════════════════════════════════════════════

    \item 제11항에 있어서,
    상기 (d) 단계는, 상기 포화 마스크의 활성 비율이 기 설정된 포화율 임계치를 초과하는 경우, 공기전도 신호의 기여를 배제하고 골전도 신호에 기반하여 음성을 복원하되, 골전도 신호의 제한된 주파수 대역을 보상하는 대역 확장을 수행하는 단계를 더 포함하는 음성 처리 방법.

    \item 제11항에 있어서,
    상기 (d) 단계는, 상기 포화 마스크가 활성인 시간-주파수 구간에서는 골전도 신호 특징의 가중치를 상대적으로 높이고, 비활성 구간에서는 공기전도 신호 특징을 우세하게 활용하는 적응적 융합을 수행하는 단계를 포함하는 음성 처리 방법.

    \item 제11항에 있어서,
    상기 음성 처리 방법은, 상기 음성 향상 모델의 추론 실패 또는 연산 오류 발생 시, 공기전도 채널 및 골전도 채널의 포화 상태에 따라 빔포밍 처리된 공기전도 신호의 우회 출력, 골전도 신호의 직접 출력, 또는 무음 처리 중 하나를 선택적으로 적용하는 폴백 단계를 더 포함하는 음성 처리 방법.

    \item 제11항에 있어서,
    상기 (e) 단계는, 좌측 및 우측 골전도 진동자를 독립 채널로 구동하여 양이간 레벨차(ILD)에 기반한 좌우 방향 구분 정보를 포함하는 스테레오 골전도 출력을 제공하는 단계를 더 포함하는 음성 처리 방법.

    \item 제11항에 있어서,
    상기 (e) 단계는, 골전도 진동자의 출력이 기 설정된 안전 임계치를 초과하지 않도록 실시간으로 제한하되, 입력 소음 수준에 따라 출력 제한 임계치를 동적으로 조절하는 단계를 더 포함하는 음성 처리 방법.

    \item 제11항에 있어서,
    상기 음성 향상 모델은, 포화, 클리핑, 및 비선형 왜곡 조건을 시뮬레이션에 의해 합성한 학습 데이터로 훈련된 것을 특징으로 하는 음성 처리 방법.

    \item 제11항에 있어서,
    상기 음성 처리 방법은, 골전도 센서와 공기전도 마이크로폰이 동시에 포화 상태인 수음 공백 구간이 검출된 경우, 직전 유효 프레임의 스펙트럼을 보간하되, 수음 공백 구간이 기 설정된 프레임 수를 초과하는 경우에는 무음 처리하는 단계를 더 포함하는 음성 처리 방법.

    \item 제11항에 있어서,
    상기 (e) 단계는, 이도 폐쇄 상태에서 발생하는 폐쇄 효과(Occlusion Effect)에 의한 저주파 과잉 증폭에 대한 보정이 적용된 골전도 진동을 출력하는 단계를 포함하는 음성 처리 방법.

    \item 제11항에 있어서,
    상기 (f) 단계는, 향상된 음성 신호를 군용 표준 오디오 인터페이스를 통해 유선으로 송신하는 단계; 및 인접한 복수의 장치 간 근거리 무선 통신 채널을 통해 직접 송수신하는 단계를 포함하는 음성 처리 방법.

\end{enumerate}


\part{요약서}

\section{요약}

\section{대표도}

\begin{figure}[H]
    \centering
    \resizebox{\textwidth}{!}{%
        \input{figures/fig1_system_block}%
    }
    \caption*{대표도 — 시스템 전체 블록도 (도 1 참조)}
    \label{fig:representative}
\end{figure}

\printbibliography
\end{document}